\part{Quantization}

\section{Geometric Quantization}

\begin{definition}[Differential forms]
A \emph{differential $k$-form} on $M$ is a smooth section of
$\Lambda^kT^*M$; write $\Omega^k(M)$ for the space of differential
$k$-forms.  The exterior
products on the cotangent spaces define
\[
\wedge:\Omega^r(M)\times\Omega^s(M)\longrightarrow\Omega^{r+s}(M).
\]
For a smooth map $F:M\to N$, define
\[
(F^*\alpha)_p(v_1,\ldots,v_k)
=\alpha_{F(p)}(dF_pv_1,\ldots,dF_pv_k).
\]
For a vector field $X$, define contraction by
\[
(\iota_X\alpha)_p(v_2,\ldots,v_k)
=\alpha_p(X_p,v_2,\ldots,v_k).
\]
The \emph{exterior derivative} is the family
$d:\Omega^k(M)\to\Omega^{k+1}(M)$ satisfying $df(X)=Xf$, the graded
Leibniz rule, $d^2=0$, and $F^*d=dF^*$.
For a vector field $X$ with local flow $\Phi_t$, define
\[
\mathcal L_X\alpha
=\left.\frac d{dt}\right|_{t=0}\Phi_t^*\alpha.
\]
\end{definition}

\begin{exercise}[Cartan calculus]
Construct $d$ locally and prove its uniqueness.  For contraction
$\iota_X$, prove Cartan's identity
\[
\mathcal L_X=d\iota_X+\iota_Xd.
\]
\end{exercise}

\begin{exercise}[Local Poincar\'e lemma]
Let $U$ be a star-shaped open neighborhood of $0$ in a finite-dimensional
real vector space.  For $\alpha\in\Omega^k(U)$ with $k>0$, define
\[
(K\alpha)_x(v_1,\ldots,v_{k-1})
=\int_0^1t^{k-1}\alpha_{tx}(x,v_1,\ldots,v_{k-1})\,dt.
\]
Use Cartan's identity for the radial flow to prove
\[
dK+Kd=\id
\]
on positive-degree forms.  Deduce that every closed form of positive degree
on $U$ is exact, with a primitive vanishing at the origin.
\end{exercise}

\begin{definition}[Cotangent phase space]
For the projection $\pi:T^*Q\to Q$, define the tautological one-form by
\[
\theta_\alpha(v)=\alpha(d\pi(v)),
\qquad \alpha\in T^*Q,\quad v\in T_\alpha(T^*Q).
\]
The canonical two-form on $T^*Q$ is $\omega=d\theta$.
\end{definition}

\begin{exercise}[From linear to cotangent phase space]
For a finite-dimensional real vector space $V$, construct the canonical
identifications
\[
TV\cong V\times V,
\qquad
T^*V\cong V\times V^*.
\]
Prove that a phase $S\in C^\infty(Q)$ determines the momentum
$p=dS_q\in T_q^*Q$ and wave covector $p/\hbar$.  For a diffeomorphism
$F:Q\to Q$, construct its cotangent lift and prove that it preserves
$\theta$ and $d\theta$.  Under $T^*V\cong V\times V^*$, prove
\[
\omega\bigl((q,p),(q',p')\bigr)=p(q')-p'(q).
\]
\end{exercise}

\begin{definition}[Symplectic manifolds]
A \emph{symplectic form} on a smooth manifold $M$ is a closed
nondegenerate two-form $\omega$.  A \emph{symplectic vector space} is a
finite-dimensional real vector space with a nondegenerate alternating
bilinear form, regarded as a constant symplectic form.  For
$f\in C^\infty(M)$, define
\[
\iota_{X_f}\omega=-df,
\qquad
\{f,g\}=-\omega(X_f,X_g).
\]
The local flow of $X_f$ is the \emph{Hamiltonian flow} of $f$.
A diffeomorphism $F:(M,\omega_M)\to(N,\omega_N)$ is a
\emph{canonical transformation} when $F^*\omega_N=\omega_M$.
\end{definition}

\begin{exercise}[Hamiltonian identities]
Prove that the canonical two-form $d\theta$ on $T^*Q$ is nondegenerate and
hence symplectic.
Prove that the Poisson bracket is antisymmetric,
satisfies the Leibniz rule and Jacobi identity, and is preserved by
canonical transformations.  For the Hamiltonian flow $\Phi_t$ of $H$,
prove
\[
\frac{d}{dt}f(\Phi_t(x))=\{f,H\}(\Phi_t(x)).
\]
\end{exercise}

\begin{exercise}[Liouville's theorem]
\label{ex:liouville-theorem}
On a $2n$-dimensional symplectic manifold, prove that Hamiltonian flow
preserves $\omega$, the volume form $\omega^n/n!$, and every density
depending only on the Hamiltonian.  On a symplectic vector space, prove
that the infinitesimal generator of a linear Hamiltonian flow has trace
zero and that the flow has determinant one.
\end{exercise}

\begin{definition}[Contact manifolds]
Let $Y$ be a smooth manifold of dimension $2n+1$.  A one-form $\alpha$ is
a \emph{contact form} when
\[
\alpha\wedge(d\alpha)^n
\]
is nowhere zero.  Its kernel $\xi=\ker\alpha$ is the associated cooriented
\emph{contact structure}; contact forms defining the same cooriented
contact structure differ by multiplication by a positive smooth function.  The
\emph{Reeb vector field} $R_\alpha$ is defined by
\[
\alpha(R_\alpha)=1,
\qquad
\iota_{R_\alpha}d\alpha=0.
\]
An immersed $n$-manifold $\iota:L\to Y$ is \emph{Legendrian} when
$\iota^*\alpha=0$.  A diffeomorphism $F:(Y,\xi)\to(Y',\xi')$ is a
\emph{contact transformation} when
$F^*\alpha'=h\alpha$ for a positive smooth function $h$.
\end{definition}

\begin{exercise}[Contactification and symplectization]
Let $(M,\omega=d\lambda)$ be an exact symplectic manifold.  On
$M\times\R$, with coordinate function $u$ on the second factor, prove that
\[
\alpha=du-\lambda
\]
is a contact form and that $R_\alpha$ is the infinitesimal generator of
translation in the $\R$ factor.  For a manifold $X$, identify
$J^1X=T^*X\times\R$ with this contactification and prove that
\[
j^1f:X\longrightarrow J^1X,
\qquad
j^1f(x)=(df_x,f(x)),
\]
is Legendrian for every $f\in C^\infty(X)$.  Prove that every section of
$J^1X\to X$ with Legendrian image is of this form.

For a contact manifold $(Y,\alpha)$ and the positive coordinate $r$ on
$\R_{>0}$, prove that
\[
\Omega=d(r\alpha)
\]
is symplectic on $\R_{>0}\times Y$.  Determine the vector field $Z$
satisfying $\iota_Z\Omega=r\alpha$ and prove that multiplication of
$\alpha$ by a positive function gives a symplectomorphic symplectization.
\end{exercise}

\begin{definition}[Integration of forms]
An \emph{orientation} of an $n$-manifold is a continuous choice of component
of the nonzero elements in $\Lambda^nT_p^*M$. The integral $\int_M\omega$ of a
compactly supported top form is defined using orientation-preserving charts
and a partition of unity.
\end{definition}

\begin{exercise}[Stokes' theorem]
Prove that the integral is independent of charts and partition of unity.
For an oriented manifold with boundary, prove
\[
\int_Md\omega=\int_{\partial M}\omega.
\]
Deduce the fundamental theorem of calculus, Green's theorem, the divergence
theorem, and the classical curl theorem. For time-dependent magnetic
two-form $B$ and electric one-form $E$, prove the equivalence of
\[
\int_{\partial S}E=-\frac d{dt}\int_SB
\qquad\text{and}\qquad
dE=-\partial_tB,
\]
recovering the observed flux--emf law.
\end{exercise}

\begin{definition}[de Rham cohomology]
A differential form $\omega$ is \emph{closed} if $d\omega=0$ and
\emph{exact} if $\omega=d\eta$ for some form $\eta$. Since $d^2=0$, the
\emph{de Rham cohomology} of $M$ is
\[
H_{\mathrm{dR}}^k(M)
=\frac{\ker(d:\Omega^k(M)\to\Omega^{k+1}(M))}
{\im(d:\Omega^{k-1}(M)\to\Omega^k(M))}.
\]
Pullback induces maps on de Rham cohomology, and the wedge product makes
$H_{\mathrm{dR}}^*(M)$ a graded-commutative algebra.
\end{definition}

\begin{exercise}[de Rham classes and periods]
Use the local Poincar\'e lemma and Stokes' theorem
to prove that an exact $k$-form has zero integral over
every closed oriented $k$-manifold mapped smoothly into $M$. Prove
$H_{\mathrm{dR}}^1(S^1)\cong\R$, with the isomorphism given by integration.

On $\R^2\setminus\{0\}$, show that
\[
\alpha=\frac{-y\,dx+x\,dy}{x^2+y^2}
\]
is closed and has integral $2\pi$ around the unit circle, hence is not exact.
Compute $H_{\mathrm{dR}}^1(\R^2\setminus\{0\})=\R\cdot[\alpha]$. Deduce that a
closed electromagnetic field strength has local potentials, while it has a
global potential $F=dA$ exactly when $[F]=0$ in de Rham cohomology.
\end{exercise}


\begin{definition}[Connections and parallel transport]
\label{def:connections-parallel-transport}
A \emph{connection} on a vector bundle $E\to M$ is an $\R$-bilinear map
$\nabla:\Gamma(TM)\times\Gamma(E)\to\Gamma(E)$, linear over smooth
functions in the first variable and satisfying
$\nabla_X(fs)=X(f)s+f\nabla_Xs$. Along a curve $\gamma$ it induces a
covariant derivative $D/dt$; its solutions $Ds/dt=0$ define
\emph{parallel transport} between endpoint fibers.
\end{definition}

\begin{definition}[Curvature and holonomy]
\label{def:curvature-holonomy}
The \emph{curvature} of $\nabla$ is the $\operatorname{End}(E)$-valued
two-form
\[
R^\nabla(X,Y)s
=\nabla_X\nabla_Ys-\nabla_Y\nabla_Xs-\nabla_{[X,Y]}s.
\]
The \emph{holonomy group} at $p$ consists of the automorphisms of $E_p$
obtained by parallel transport around loops based at $p$.
\end{definition}

\begin{exercise}[Parallel transport and curvature]
Prove existence, uniqueness, composition, and inverse properties of parallel
transport. Prove that curvature is tensorial and gives the leading failure
of transport around an infinitesimal parallelogram to return a vector to
itself. Deduce that a flat connection has path-independent transport on every
simply connected coordinate domain.
\end{exercise}

\begin{definition}[Principal connections]
For a Lie group $G$, a \emph{principal $G$-bundle} is a smooth map
$P\to M$ with a free right $G$-action, transitive on each fiber, and local
$G$-equivariant trivializations.  For a representation
$G\to\operatorname{GL}(W)$, its
associated vector bundle is
$P\times_GW=(P\times W)/((pg,w)\sim(p,gw))$.
A \emph{principal connection} is an
equivariant $\mathfrak g$-valued one-form $A$ on $P$ reproducing fundamental
vertical fields. It induces connections on every associated bundle, and its
\emph{curvature} is
$F_A=dA+\frac12[A\wedge A]$. A \emph{gauge transformation} is a
$G$-equivariant bundle automorphism covering $\id_M$.
\end{definition}

\begin{exercise}[Gauge covariance and Bianchi identity]
Prove $d_AF_A=0$.  For a gauge transformation $\Phi$, prove
$F_{\Phi^*A}=\Phi^*F_A$ and that holonomy changes by conjugacy.  Prove that
a flat connection has trivial holonomy on contractible loops and is
gauge-equivalent to a product connection on a connected base exactly when
its holonomy representation is trivial.  For $G=U(1)$, compute the holonomy
around a
generator of the fundamental group and exhibit the obstruction to removing
a flat connection globally.
\end{exercise}



\begin{definition}[Prequantum line bundles]
A \emph{prequantum line bundle} over $(M,\omega)$ is a Hermitian line bundle
$L\to M$ with a unitary connection $\nabla$ satisfying
\[
F_\nabla=-\frac{i}{\hbar}\omega.
\]
For $f\in C^\infty(M)$, its \emph{Kostant--Souriau operator} is
\[
\mathcal Q_{\rm pre}(f)=-i\hbar\nabla_{X_f}+f.
\]
\end{definition}

\begin{definition}[Metaplectic correction]
A \emph{complex polarization} is an involutive Lagrangian subbundle
$P\subset T_\C M$.  Its annihilator and canonical line are
\[
P^\circ=\{\alpha\in T_\C^*M:\alpha|_P=0\},
\qquad
K_P=\Lambda^{\rm top}P^\circ.
\]
A \emph{half-form bundle} is a line bundle $\delta_P$ with
$\delta_P^{\otimes2}\cong K_P$, equipped with the partial connection along
$P$ induced by Lie derivatives.  A corrected polarized section of
$L\otimes\delta_P$ is covariantly constant along $P$ and square-integrable
for the density determined by the half-form pairing.  The corrected quantum
space is its Hilbert completion when this pairing is positive.
\end{definition}

\begin{exercise}[Prequantum integrality]
Prove that a prequantum line bundle exists if and only if the period
condition
\[
\frac1{2\pi\hbar}\int_\Sigma f^*\omega\in\Z
\]
holds for every smooth map from a closed oriented surface
$f:\Sigma\to M$.  Construct the line-bundle transition functions from
local symplectic potentials and prove that the period condition is exactly
their cocycle condition.
\end{exercise}

\begin{exercise}[Kostant--Souriau commutators]
Prove
\[
[\mathcal Q_{\rm pre}(f),\mathcal Q_{\rm pre}(g)]
=i\hbar\mathcal Q_{\rm pre}(\{f,g\}).
\]
\end{exercise}

\begin{exercise}[Contactification]
Let $\pi:S(L)\to M$ be the unit-circle bundle of $L$, let
$A\in\Omega^1(S(L);i\R)$ be its connection form normalized by
$A(\xi)=i$ on the circle generator $\xi$, and put $\alpha=-iA$.
Prove
\[
d\alpha=-\frac1\hbar\pi^*\omega,
\]
that $\alpha$ is a contact form with Reeb vector field $\xi$, and that
sections of $L$ correspond to functions $F:S(L)\to\C$ satisfying
$F(ue^{i\theta})=e^{-i\theta}F(u)$.
\end{exercise}

\begin{exercise}[Polarized quantum spaces]
For $T^*V$ with the vertical polarization, recover $L^2(V)$ and the
operators
\[
\psi(x)\longmapsto \ell(x)\psi(x),
\qquad
\psi(x)\longmapsto-i\hbar\,d\psi_x(v).
\]
For a compact
K\"ahler manifold with $P=T^{0,1}M$, prove $P^\circ=T^{*1,0}M$ and
$K_P=K_M$.  For a positive prequantum line bundle, identify the
corrected quantum space with
$H^0(M,L\otimes K_M^{1/2})$.
\end{exercise}


\begin{definition}[Theta functions with characteristics]
For $a,b\in\R$, $z\in\C$, and $\tau$ in the upper half-plane, define
\[
\vartheta\!\begin{bmatrix}a\\b\end{bmatrix}(z\mid\tau)
=\sum_{n\in\Z}
\exp\!\left(
\pi i(n+a)^2\tau+2\pi i(n+a)(z+b)
\right).
\]
\end{definition}

\begin{exercise}[Landau theta functions]
Let
$Q=\R^2/(L_x\Z\oplus L_y\Z)$ have area $A=L_xL_y$, let a particle of
charge $e$ experience the constant magnetic field $B\,dx\wedge dy$, and
assume $eB>0$.  Put
\[
N=\frac{eBA}{2\pi\hbar},
\qquad
\tau=i\frac{L_y}{L_x},
\qquad
z=\frac{x}{L_x}+\tau\frac{y}{L_y}.
\]
Use the prequantization condition to prove that a magnetic line bundle
$L\to Q$ exists exactly when $N\in\Z$.  For $N>0$, use the holomorphic-gauge
automorphy laws
\[
f(z+1)=f(z),
\qquad
f(z+\tau)=e^{-\pi iN\tau-2\pi iNz}f(z)
\]
to prove the magnetic-translation decomposition
\[
H^0(Q,L)
=\bigoplus_{r\in\Z/N\Z}
\C\,\vartheta\!\begin{bmatrix}r/N\\0\end{bmatrix}(Nz\mid N\tau).
\]
Derive the automorphy and magnetic-translation eigenvalue laws directly
from the defining series.  Relate the holomorphic-gauge sections to
physical wavefunctions by their common unitary-gauge factor.

For magnetic translations by $L_x/N$ and $L_y/N$, use holonomy around a
fundamental magnetic cell to prove
\[
UV=e^{2\pi i/N}VU.
\]
For the covariant Landau Hamiltonian, construct the raising and lowering
operators and prove
\[
E_n=\hbar\frac{eB}{m}\left(n+\frac12\right),
\qquad
\dim E_{\widehat H}(\{E_n\})H=N.
\]
Deduce the cyclotron resonance spacing, the degeneracy density
$eB/(2\pi\hbar)$, and the magnetic-translation interference phase.
State how measurements of these three quantities separately test the
energy, flux-integrality, and theta-function predictions.
\end{exercise}

\section{Deformation Quantization}

\begin{definition}[Coadjoint action and stabilizers]
For a Lie group $G$ with Lie algebra $\mathfrak g$, define
\[
\Ad_g=d(h\mapsto ghg^{-1})_e,
\qquad
(\Ad_g^*\ell)(X)=\ell(\Ad_{g^{-1}}X).
\]
For $\ell\in\mathfrak g^*$, define
\[
\mathcal O_\ell=\Ad^*(G)\ell,
\qquad
G_\ell=\{g\in G:\Ad_g^*\ell=\ell\},
\]
and
\[
\mathfrak g_\ell
=\{X\in\mathfrak g:\ell([X,Y])=0
\text{ for every }Y\in\mathfrak g\}.
\]
For $X\in\mathfrak g$, set
\[
(\ad_X^*\ell)(Y)=-\ell([X,Y]).
\]
\end{definition}

\begin{definition}[Kirillov form]
On $T_\ell\mathcal O_\ell$, define
\[
\omega_\ell(\ad_X^*\ell,\ad_Y^*\ell)=\ell([X,Y]).
\]
For $\eta=\Ad_g^*\ell$, transport $\omega_\ell$ by the coadjoint
action to define $\omega_\eta$.
\end{definition}

\begin{exercise}[Geometry of coadjoint orbits]
Prove
\[
T_\ell\mathcal O_\ell
=\{\ad_X^*\ell:X\in\mathfrak g\}
\cong\mathfrak g/\mathfrak g_\ell.
\]
Prove that the Kirillov form is well-defined, alternating, nondegenerate,
and $G$-invariant.  Prove the cyclic identity
\[
\ell([X,[Y,Z]])+\ell([Y,[Z,X]])+\ell([Z,[X,Y]])=0
\]
and identify it with the closedness identity for the invariant two-form
$\omega$ on $\mathcal O_\ell$.
\end{exercise}


\begin{definition}[Projective representations]
A \emph{normalized multiplier} on a locally compact group $G$ is a
continuous map
$\sigma:G\times G\to\T$ satisfying
\[
\sigma(x,e)=\sigma(e,x)=1,
\qquad
\sigma(x,y)\sigma(xy,z)=\sigma(y,z)\sigma(x,yz).
\]
A \emph{$\sigma$-projective unitary representation} is a strongly continuous map
$U:G\to U(H)$ satisfying
\[
U_xU_y=\sigma(x,y)U_{xy},
\qquad U_e=I.
\]
Multipliers $\sigma$ and $\sigma'$ are \emph{cohomologous} when there is a
continuous map $b:G\to\T$, $b(e)=1$, such that
\[
\sigma'(x,y)=b(x)b(y)\overline{b(xy)}\,\sigma(x,y).
\]
\end{definition}


\begin{definition}[Weyl multipliers]
Let $(E,\omega)$ be a symplectic vector space.
For $\lambda\in\R$, define
\[
\sigma_{\lambda\omega}(z,z')=e^{i\lambda\omega(z,z')/2}.
\]
For physical momentum variables and $\hbar>0$, the \emph{Weyl multiplier}
is $\sigma_{\omega/\hbar}$.
\end{definition}

\begin{exercise}[The Weyl multiplier]
Prove that $\sigma_{\lambda\omega}$ is a normalized multiplier for every
$\lambda\in\R$.
\end{exercise}


\begin{definition}[Commutator bicharacter]
For a multiplier $\sigma$ on an abelian group $A$, define
\[
c_\sigma(x,y)=\sigma(x,y)\overline{\sigma(y,x)}.
\]
\end{definition}

\begin{exercise}[Commutator bicharacter]
Prove that $c_\sigma$ is an alternating bicharacter and depends only on the
cohomology class of $\sigma$.  For the Weyl multiplier on $V\oplus V^*$,
prove
\[
c_{\sigma_{\omega/\hbar}}\bigl((q,p),(q',p')\bigr)
=e^{i(p(q')-p'(q))/\hbar}.
\]
Prove
\[
c_{\sigma_{\lambda\omega}}(z,z')=e^{i\lambda\omega(z,z')}.
\]
\end{exercise}


\begin{definition}[Heisenberg group]
For a symplectic vector space $(E,\omega)$, its \emph{Heisenberg group} is
\[
\mathbb H(E,\omega)=E\times\R
\]
with multiplication
\[
(z,s)(z',s')
=(z+z',s+s'+\tfrac12\omega(z,z')).
\]
Write
$\mathbb H_n=\mathbb H(\R^n\oplus(\R^n)^*,\omega)$ for the canonical
symplectic form on $\R^n\oplus(\R^n)^*$.  Its Lie algebra is
$\mathfrak h(E,\omega)=E\oplus\R Z$ with
\[
[(z,sZ),(z',s'Z)]=\omega(z,z')Z.
\]
For $E=V\oplus V^*$ and $\lambda\ne0$, the \emph{Schr\"odinger
representation} on $L^2(V)$ is
\[
\bigl(\pi_\lambda(q,p,s)\xi\bigr)(x)
=e^{i\lambda(s+p(x-q/2))}\xi(x-q).
\]
For $\hbar>0$, $\ell\in V^*$, and $v\in V$, define on $\mathcal S(V)$
\[
Q_\ell\psi(x)=\ell(x)\psi(x),
\qquad
P_v\psi(x)=-i\hbar\,d\psi_x(v).
\]
\end{definition}

\begin{definition}[Polarization]
For a connected simply connected nilpotent group $G$ and
$\ell\in\mathfrak g^*$, a \emph{polarization} is a Lie subalgebra
$\mathfrak p\subseteq\mathfrak g$ satisfying
\[
\ell([\mathfrak p,\mathfrak p])=0,
\qquad
\dim\mathfrak p
=\frac{\dim\mathfrak g+\dim\mathfrak g_\ell}{2}.
\]
For $P=\exp(\mathfrak p)$, set
$\chi_\ell(\exp X)=e^{i\ell(X)}$.  The induced representation
$\operatorname{Ind}_P^G\chi_\ell$ acts by left translation on the
$L^2$-sections over $G/P$ satisfying
$F(gp)=\chi_\ell(p)^{-1}F(g)$.
\end{definition}

\begin{exercise}[Central extensions]
Prove that $\mathbb H(E,\omega)$ is a locally compact group with center
$\{0\}\times\R$ and that
\[
\pi_\lambda(0,s)=e^{i\lambda s}I.
\]
Prove that projective representations of $E$ with multiplier
$e^{i\lambda\omega/2}$ correspond to representations of
$\mathbb H(E,\omega)$ with central character $e^{i\lambda s}$.
Let $p:\mathbb H(E,\omega)\to E$ be the quotient map and define
\[
\alpha_{(z,s)}(v,r)=r-\frac12\omega(z,v).
\]
Prove that $\alpha$ is left-invariant,
\[
d\alpha=-p^*\omega,
\]
and that $\alpha$ is a contact form whose Reeb flow is translation in the
central direction.
\end{exercise}

\begin{exercise}[Canonical commutation and Weyl relations]
Prove
\[
[Q_\ell,P_v]=i\hbar\ell(v)I.
\]
For
\[
(T_q\psi)(x)=\psi(x-q),
\qquad
(M_p^\hbar\psi)(x)=e^{ip(x)/\hbar}\psi(x),
\]
prove
\[
M_p^\hbar T_q=e^{ip(q)/\hbar}T_qM_p^\hbar.
\]
For
\[
W_\hbar(q,p)=e^{-ip(q)/(2\hbar)}M_p^\hbar T_q,
\]
prove
\[
W_\hbar(z)W_\hbar(z')
=e^{i\omega(z,z')/(2\hbar)}W_\hbar(z+z')
\]
and identify $W_\hbar(z)$ with
$\pi_{1/\hbar}(z,0)$ under the canonical set-theoretic section
$z\mapsto(z,0)$.  Prove that this section is not a homomorphism and that
its failure to preserve multiplication is exactly the displayed multiplier.
\end{exercise}


\begin{exercise}[Schr\"odinger model]
Identify
\[
\mathfrak h(E,\omega)^*=E^*\oplus\R Z^*.
\]
For $(\xi,\lambda)\in E^*\oplus\R Z^*$, prove
\[
\Ad_{(z,s)}^*(\xi,\lambda)
=\bigl(\xi-\lambda\omega(z,{\cdot}),\lambda\bigr).
\]
Deduce that the orbits with $\lambda=0$ are points and that, for
$\lambda\ne0$,
\[
\mathcal O_\lambda
=\{(\xi,\lambda):\xi\in E^*\}.
\]
Prove that
\[
E\longrightarrow\mathcal O_\lambda,
\qquad
z\longmapsto\bigl(-\lambda\omega(z,{\cdot}),\lambda\bigr),
\]
pulls the Kirillov form back to $\lambda\omega$.

For $E=V\oplus V^*$ and
$\mathfrak p=V^*\oplus\R Z$, prove that $\mathfrak p$ is a polarization
at $(0,\lambda)$.  Identify
$\mathbb H(E,\omega)/\exp(\mathfrak p)$ with $V$ and compute the
induced action to obtain
\[
\bigl(\pi_\lambda(q,p,s)\psi\bigr)(x)
=e^{i\lambda(s+p(x-q/2))}\psi(x-q).
\]
For $\lambda=0$, recover every character of $E$ from a point orbit.
\end{exercise}


\begin{exercise}[Irreducibility of the Schr\"odinger representation]
Prove that an operator on $L^2(V)$ commuting with every translation and
every modulation is scalar: first use the Fourier transform to identify the
commutant of translations, and then impose commutation with modulations.
Deduce that $\pi_\lambda$ is irreducible for $\lambda\ne0$.
\end{exercise}

\begin{exercise}[Stone--von Neumann theorem]
Let $U$ be a strongly continuous irreducible unitary representation of
$\mathbb H(V\oplus V^*,\omega)$ with central character
$U(0,s)=e^{i\lambda s}I$, where $\lambda\ne0$.  Integrate Schwartz
functions on $V\oplus V^*$ against the canonical section of $U$.  Use
Fourier inversion and a Gaussian approximate identity to show that this
integrated representation contains the same nonzero finite-rank operators
as the Schr\"odinger model.  Construct a nonzero intertwiner and use
irreducibility to prove $U\simeq\pi_\lambda$.
\end{exercise}

\begin{definition}[Trace and Hilbert--Schmidt classes]
For $x,y\in H$, let
\[
\theta_{x,y}(\xi)=\langle y,\xi\rangle x.
\]
The \emph{trace class} $S_1(H)$ consists of operators admitting a representation
\[
A=\sum_{j=1}^{\infty}\theta_{x_j,y_j},
\qquad
\sum_{j=1}^{\infty}\|x_j\|\,\|y_j\|<\infty,
\]
where the series converges in operator norm, with norm the infimum of the
displayed sums.  Define
\[
\Tr(A)=\sum_{j=1}^{\infty}\langle y_j,x_j\rangle.
\]
The \emph{Hilbert--Schmidt class} is
\[
S_2(H)=\{A\in B(H):A^*A\in S_1(H)\},
\qquad
\|A\|_2^2=\Tr(A^*A).
\]
\end{definition}

\begin{exercise}[Trace ideals]
Prove that $S_1(H)$, $\Tr$, and $S_2(H)$ are well-defined.  Prove that
$\langle A,B\rangle_{S_2}=\Tr(A^*B)$ makes $S_2(H)$ a Hilbert space and
that the pairing $(A,T)\mapsto\Tr(AT)$ identifies $S_1(H)^*$ isometrically
with $B(H)$.
\end{exercise}


\begin{definition}[Coefficient transforms]
Let $(X,\mu)$ be a measure space and let $U:X\to U(H)$ be a weakly
measurable map.  The \emph{coefficient transform} of the unitary system
$U$ is
\[
\mathcal V_U:S_1(H)\longrightarrow L^\infty(X),
\qquad
\mathcal V_U(A)(x)=\Tr(AU_x).
\]
\end{definition}

\begin{definition}[Fourier--Wigner transform]
Let $V$ be an $n$-dimensional real vector space with Lebesgue measure $dq$,
let $\hbar>0$, and equip $V^*$ with the dual measure $dp$ for which
$(2\pi\hbar)^{-n/2}\int_Ve^{-ip(q)/\hbar}f(q)\,dq$ is unitary on $L^2$.
For $q\in V$ and $p\in V^*$, define
\[
(T_q\xi)(x)=\xi(x-q),
\qquad
(M_p^\hbar\xi)(x)=e^{ip(x)/\hbar}\xi(x).
\]
For $z=(q,p)\in V\oplus V^*$, define
\[
W_\hbar(z)=e^{-ip(q)/(2\hbar)}M_p^\hbar T_q.
\]
The \emph{Fourier--Wigner transform} is
\[
\mathcal W_\hbar:S_1(L^2(V))\longrightarrow C_b(V\oplus V^*),
\qquad
\mathcal W_\hbar(A)(z)=(2\pi\hbar)^{-n/2}\Tr(AW_\hbar(z)).
\]
\end{definition}

\begin{exercise}[Fourier--Wigner Plancherel]
Prove
\[
W_\hbar(z)^*=W_\hbar(-z).
\]
For $A,B\in S_1(L^2(V))\cap S_2(L^2(V))$, prove the Fourier--Wigner
orthogonality identity
\[
\int_{V\oplus V^*}
\overline{\mathcal W_\hbar(A)(z)}\mathcal W_\hbar(B)(z)\,dz
=\Tr(A^*B).
\]
Deduce that $\mathcal W_\hbar$ extends uniquely to a unitary map
\[
S_2(L^2(V))\longrightarrow L^2(V\oplus V^*)
\]
and, whenever $\mathcal W_\hbar(A)\in L^1(V\oplus V^*)$, prove
\[
A=(2\pi\hbar)^{-n/2}\int_{V\oplus V^*}
\mathcal W_\hbar(A)(z)W_\hbar(z)^*\,dz
\]
in the weak operator topology.
\end{exercise}


\begin{exercise}[Schr\"odinger character]
For $V=\R^n$, let $f\in C^\infty(\mathbb H_n)$ satisfy
\[
(q,p,s)\longmapsto f\bigl(\exp(q,p,s)\bigr)
\in\mathcal S(\R^{2n+1})
\]
under the canonical identification
$\mathfrak h_n=\R^n\oplus(\R^n)^*\oplus\R$.  With Haar measure
$dq\,dp\,ds$, put
\[
\pi_\lambda(f)
=\int_{\mathbb H_n}f(g)\pi_\lambda(g)\,dg.
\]
Prove
\[
\Tr\pi_\lambda(f)
=(2\pi)^n|\lambda|^{-n}
\int_\R f(0,0,s)e^{i\lambda s}\,ds.
\]
Derive this identity directly from the Schr\"odinger kernel and identify
the factor $|\lambda|^{-n}$ with the symplectic-volume scaling of
$\mathcal O_\lambda$.
\end{exercise}

\begin{exercise}[Heisenberg Plancherel and inversion]
\label{ex:heisenberg-harmonic-analysis}
For $f\in L^1(\mathbb H_n)\cap L^2(\mathbb H_n)$, use Haar measure
$dq\,dp\,ds$ and the scalar Fourier transform
$\widehat h(\lambda)=\int_\R h(s)e^{-i\lambda s}\,ds$, and define
\[
\widehat f(\lambda)
=\int_{\mathbb H_n}f(g)\pi_\lambda(g)^*\,dg.
\]
Compute the integral kernel of $\widehat f(\lambda)$ and apply Euclidean
Plancherel in the central and phase-space variables.  Deduce the Plancherel measure
$(2\pi)^{-(n+1)}|\lambda|^n\,d\lambda$ and prove
\[
\|f\|_2^2
=(2\pi)^{-(n+1)}\int_{\R\setminus\{0\}}
\|\widehat f(\lambda)\|_{\rm HS}^2|\lambda|^n\,d\lambda.
\]
Prove the inversion formula and the distributional orthogonality of Weyl
operators.  Prove that representations with trivial central character have
zero Plancherel measure.
\end{exercise}

\begin{definition}[Weyl quantization]
For $a\in\mathcal S(V\oplus V^*)$, define
$\operatorname{Op}_\hbar^W(a):\mathcal S(V)\to\mathcal S(V)$ by
\[
\bigl(\operatorname{Op}_\hbar^W(a)\psi\bigr)(x)
=\frac{1}{(2\pi\hbar)^n}
\int_{V\times V^*}
e^{ip(x-y)/\hbar}
a\!\left(\frac{x+y}{2},p\right)\psi(y)\,dy\,dp.
\]
For polynomial symbols, use the same formula as an oscillatory integral on
$\mathcal S(V)$.
\end{definition}

\begin{definition}[Wigner distributions]
For $\psi,\varphi\in\mathcal S(V)$, define
\[
W_\hbar(\psi,\varphi)(q,p)
=\frac{1}{(2\pi\hbar)^n}
\int_Ve^{-ip(y)/\hbar}
\overline{\psi(q-y/2)}\,\varphi(q+y/2)\,dy.
\]
Write $W_\hbar[\psi]=W_\hbar(\psi,\psi)$.
\end{definition}

\begin{exercise}[Weyl calculus and the Moyal product]
For $\ell\in V^*$ and $v\in V$, prove
\[
\operatorname{Op}_\hbar^W\bigl((q,p)\mapsto\ell(q)\bigr)=Q_\ell,
\qquad
\operatorname{Op}_\hbar^W\bigl((q,p)\mapsto p(v)\bigr)=P_v.
\]
Prove
\[
\operatorname{Op}_\hbar^W(a)^*
=\operatorname{Op}_\hbar^W(\overline a),
\qquad
\|\operatorname{Op}_\hbar^W(a)\|_{\rm HS}
=(2\pi\hbar)^{-n/2}\|a\|_2,
\]
and
\[
\langle\psi,\operatorname{Op}_\hbar^W(a)\varphi\rangle
=\int_{V\oplus V^*}a(q,p)W_\hbar(\psi,\varphi)(q,p)\,dq\,dp.
\]
Prove that for $a,b\in\mathcal S(V\oplus V^*)$ there is a unique
$a\star_\hbar b\in\mathcal S(V\oplus V^*)$ satisfying
\[
\operatorname{Op}_\hbar^W(a\star_\hbar b)
=\operatorname{Op}_\hbar^W(a)\operatorname{Op}_\hbar^W(b).
\]
Derive the oscillatory-integral formula for $\star_\hbar$ and prove
\[
a\star_\hbar b
=ab+\frac{i\hbar}{2}\{a,b\}+O(\hbar^2),
\qquad
\frac{a\star_\hbar b-b\star_\hbar a}{i\hbar}
=\{a,b\}+O(\hbar^2)
\]
in the Schwartz-symbol topology.
\end{exercise}

\begin{exercise}[Groenewold--van Hove obstruction]
Specialize to $V=\R$ and write $Q=Q_{\id_\R}$ and $P=P_1$.
Prove that the polynomials of degree at most two form a Lie subalgebra for
the Poisson bracket and that Weyl quantization is an exact Lie-algebra
quantization on this subalgebra.  Quantize $q^2p^2$ using
\[
q^2p^2=\frac19\{q^3,p^3\}
=\frac13\{q^2p,qp^2\}
\]
and prove that there is no linear quantization of the polynomial Poisson
algebra satisfying
\[
\mathcal Q(1)=I,\qquad
\mathcal Q(q)=Q,\qquad
\mathcal Q(p)=P,\qquad
\mathcal Q(q^m)=Q^m,\qquad
\mathcal Q(p^m)=P^m,
\]
together with
\[
\mathcal Q(\{f,g\})=(i\hbar)^{-1}
[\mathcal Q(f),\mathcal Q(g)]
\]
on a common invariant domain.
\end{exercise}


\begin{definition}[Classical and quantum evolution]
Let $H\in C^\infty(V\oplus V^*,\R)$ have bounded derivatives of order at least
two, let $\Phi_t$ be its Hamiltonian flow, and let
$\operatorname{Op}_\hbar^W(H)$ denote its self-adjoint realization.  Define
\[
U_\hbar(t)=e^{-it\operatorname{Op}_\hbar^W(H)/\hbar}.
\]
For $a\in\mathcal S(V\oplus V^*)$, define
\[
a_t=a\circ\Phi_t,
\qquad
A_\hbar(t)=U_\hbar(t)^*\operatorname{Op}_\hbar^W(a)U_\hbar(t).
\]
\end{definition}

\begin{exercise}[Heisenberg equation and Egorov's theorem]
Prove
\[
\frac{d}{dt}a_t=\{a_t,H\},
\qquad
\frac{d}{dt}A_\hbar(t)
=\frac{i}{\hbar}
[\operatorname{Op}_\hbar^W(H),A_\hbar(t)].
\]
For every $T>0$, prove Egorov's theorem
\[
\sup_{|t|\leq T}
\left\|
A_\hbar(t)-\operatorname{Op}_\hbar^W(a\circ\Phi_t)
\right\|=O_T(\hbar).
\]
Prove exact equality for quadratic $H$.  Derive Ehrenfest's equations for
the expectations of $Q_\ell$ and $P_v$ for every $\ell\in V^*$ and
$v\in V$.
\end{exercise}

\begin{exercise}[Linear symplectic covariance]
Let $S$ be a linear symplectic automorphism of $V\oplus V^*$.  Construct
unitary implementers for cotangent lifts of linear automorphisms of $V$,
for symplectic shears defined by quadratic forms, and for the exchange of
position and momentum implemented by the Fourier transform.  Prove for
these generators that
\[
\mu(S)\operatorname{Op}_\hbar^W(a)\mu(S)^*
=\operatorname{Op}_\hbar^W(a\circ S^{-1}),
\qquad
W_\hbar[\mu(S)\psi]=W_\hbar[\psi]\circ S^{-1}.
\]
Deduce the covariance for their products and prove that each implementer is
determined up to phase.  Show that coherent choices form the metaplectic
double cover, whose two lifts of a fixed symplectic map differ by sign.  Use
the Groenewold--van Hove obstruction to show why exact covariance cannot
extend to every nonlinear canonical transformation.
\end{exercise}

\begin{exercise}[Wigner--Laguerre identity]
Use the generalized Laguerre polynomials and their orthogonality from
Determinantal Ensembles.  For
\[
H_{\rm cl}(q,p)=\frac{p^2}{2m}+\frac{m\omega^2q^2}{2},
\qquad
\psi_n(q)=\left(\frac{m\omega}{\hbar}\right)^{1/4}
h_n\!\left(\sqrt{\frac{m\omega}{\hbar}}q\right),
\]
prove
\[
W_\hbar[\psi_n](q,p)
=\frac{(-1)^n}{\pi\hbar}
e^{-2H_{\rm cl}(q,p)/(\hbar\omega)}
L_n\!\left(\frac{4H_{\rm cl}(q,p)}{\hbar\omega}\right).
\]
Verify the position and momentum marginals and identify the zeros and sign
changes of the Wigner distribution from those of $L_n$.
\end{exercise}


\begin{exercise}[Laguerre tomography]
For the dimensionless quadrature
\[
X_\theta=\sqrt{\frac{m\omega}{\hbar}}Q\cos\theta
+\frac{P}{\sqrt{m\hbar\omega}}\sin\theta,
\]
prove that its probability density in a state $\psi$ is the Radon
transform of $W_\hbar[\psi]$ along the corresponding phase-space lines.
Derive the Fourier-slice inversion formula reconstructing the Wigner
distribution from the family of quadrature densities.  For the oscillator
state $\psi_n$, use the Wigner--Laguerre identity to prove that the
reconstructed radial profile is
\[
\frac{(-1)^n}{\pi\hbar}e^{-2H_{\rm cl}/(\hbar\omega)}
L_n\!\left(\frac{4H_{\rm cl}}{\hbar\omega}\right),
\]
while every quadrature density is a scaled
$e^{-x^2}H_n(x)^2$.  Deduce the predicted number of quadrature nodes, the
radial zero circles determined by $L_n$, and the sign at the phase-space
origin.  State how homodyne quadrature data distinguish adjacent number
states and detect Wigner negativity after inverse Radon reconstruction.
\end{exercise}

\begin{definition}[Poisson manifolds]
A \emph{Poisson manifold} is a smooth manifold $M$ with a bracket on
$C^\infty(M)$ that is a Lie bracket and a derivation in each argument.
Equivalently, there is a bivector $\Pi\in\Gamma(\Lambda^2TM)$ satisfying
\[
\{f,g\}=\Pi(df,dg),
\qquad [\Pi,\Pi]_{\rm Sch}=0.
\]
\end{definition}

\begin{definition}[Star products]
A \emph{differential star product} on a Poisson manifold is an associative
$\C[[\hbar]]$-bilinear product on $C^\infty(M)[[\hbar]]$ of the form
\[
f\star g=fg+\sum_{r\geq1}\hbar^rC_r(f,g),
\]
where the $C_r$ are bidifferential operators, $1$ is a unit, and
\[
C_1(f,g)-C_1(g,f)=i\{f,g\}.
\]
Two star products are \emph{equivalent} when they are intertwined by an
operator $T=\id+\sum_{r\geq1}\hbar^rT_r$ with differential $T_r$.
\end{definition}

\begin{exercise}[Affine Moyal product]
Let $(E,\omega)$ be a symplectic vector space and let
$\Pi:E^*\to E$ be the inverse of $v\mapsto\omega(v,-)$, and write
$\Pi(k,\ell)=\ell(\Pi(k))$.  Let $D_\Pi$ be the constant bidifferential
operator characterized on decomposable functions by
\[
D_\Pi(f\otimes g)(x,y)
=d_yg\bigl(\Pi(d_xf)\bigr).
\]
On $C^\infty(E)[[\hbar]]$, define
\[
f\star g
=m\circ\exp\!\left(\frac{i\hbar}{2}D_\Pi\right)(f\otimes g),
\]
where $m$ restricts a function on $E\times E$ to the diagonal.  Prove
associativity from the commutation of the induced constant bidifferential
operators on $E^3$, and prove
\[
f\star g=fg+\frac{i\hbar}{2}\{f,g\}+O(\hbar^2).
\]
For $k,\ell\in E^*$ put $e_k(x)=e^{ik(x)}$ and derive
\[
e_k\star e_\ell
=e^{-i\hbar\Pi(k,\ell)/2}e_{k+\ell}.
\]
Prove associativity directly by checking the multiplier identity for the
phase.  Differentiate in $k$ and $\ell$ to recover the canonical
commutation relations.  Prove that every affine symplectic transformation
of $E$ preserves $\star$.  For $E=V\oplus V^*$, prove that the restriction
to Schwartz functions is the formal $\hbar$-expansion of the Weyl product
defined by operator composition above.
\end{exercise}

\begin{definition}[Magnetic lattice Hamiltonian]
Let $U,V$ be unitary translations satisfying
\[
UV=e^{2\pi i\alpha}VU,
\]
where $\alpha$ is the magnetic flux through one lattice cell in units of
the flux quantum.  The magnetic nearest-neighbor Hamiltonian is
\[
H_\alpha=U+U^*+V+V^*.
\]
In the position realization on $\ell^2(\Z)$, with transverse phase
$\theta$, it is the Harper operator
\[
(H_{\alpha,\theta}\psi)_m
=\psi_{m+1}+\psi_{m-1}
+2\cos(2\pi\alpha m+\theta)\psi_m.
\]
\end{definition}

\begin{exercise}[Chebyshev magnetic bands]
Derive the magnetic-translation relation from the Moyal plane-wave
identity.  For rational flux $\alpha=p/q$, prove that the Harper equation
has $q$-periodic coefficients and reduce it by the Floquet condition
$\psi_{m+q}=e^{iqk}\psi_m$ to a $q$-dimensional spectral problem.
At zero flux, use the Chebyshev identity
\[
T_q(\cos t)=\cos(qt)
\]
to derive the $q$-step discriminant
\[
\Delta_{0,q}(E,\theta)
=2T_q\!\left(\frac{E-2\cos\theta}{2}\right).
\]
At half a flux quantum, compute the two bands
\[
E_\pm(k,\theta)
=\pm2\sqrt{\cos^2k+\cos^2\theta}.
\]
Deduce the flux-dependent band edges and the touching energy, and state
how microwave, photonic, or electronic lattice spectroscopy tests the
Chebyshev discriminant and the half-flux splitting.
Prove that the commutator phase $e^{2\pi i\alpha}$ is the phase accumulated
around one lattice cell.  Show that nonintegral total flux on a periodic
sample obstructs consistent torus boundary phases, while integral total
flux makes their corner cocycle trivial.  Prove that nonzero integral flux
still obstructs a single globally periodic vector potential.
\end{exercise}



\section{Quantum Groups}



\begin{definition}[Poisson--Lie groups]
A \emph{Poisson--Lie group} is a Lie group $G$ with a Poisson structure for
which multiplication $G\times G\to G$ is Poisson.  A \emph{Lie bialgebra}
is a Lie algebra $\mathfrak g$ with a linear map
\[
\delta:\mathfrak g\to\mathfrak g\wedge\mathfrak g
\]
that is a $1$-cocycle for the adjoint action and whose transpose defines a
Lie bracket on $\mathfrak g^*$.  It is \emph{coboundary} when
\[
\delta(x)=[x\otimes1+1\otimes x,r]
\]
for some $r\in\mathfrak g\otimes\mathfrak g$.
\end{definition}

\begin{exercise}[Infinitesimal Poisson--Lie structure]
Differentiate a Poisson--Lie structure at the identity and prove that it
defines a Lie bialgebra.  For a coboundary Lie bialgebra, define
\[
\operatorname{CYB}(r)
=[r_{12},r_{13}]+[r_{12},r_{23}]+[r_{13},r_{23}].
\]
Prove that the co-Jacobi identity is equivalent to
$\operatorname{CYB}(r)$ being invariant under the diagonal adjoint action.
Under the additional triangular hypothesis
$\operatorname{CYB}(r)=0$, recover the classical Yang--Baxter equation;
state the modified classical Yang--Baxter equation as
$\operatorname{CYB}(r)=\Omega$ for an invariant tensor $\Omega$ and
distinguish the case $\Omega\ne0$.
\end{exercise}


\begin{definition}[Bialgebras and Hopf algebras]
A \emph{bialgebra} is a unital associative algebra $H$ with unital algebra
homomorphisms
\[
\Delta:H\to H\otimes H,
\qquad
\varepsilon:H\to\C
\]
satisfying coassociativity and the counit identities.  A \emph{Hopf
algebra} is a bialgebra with a linear antipode $S:H\to H$ satisfying
\[
m(S\otimes\id)\Delta
=\eta\varepsilon
=m(\id\otimes S)\Delta.
\]
\end{definition}

\begin{definition}[Deformation parameters]
The generic quantum-group parameter is $0<v<1$, the basic-hypergeometric
base is $Q$, the root-of-unity parameter is $\zeta$, the DAHA parameters
are $q,t$, the elliptic nome is $p$, and $\kappa\ne0$ is the independent
$q$KZ shift.
\end{definition}

\begin{definition}[Compact real form of \texorpdfstring{$U_v(\mathfrak{sl}_2)$}{Uv(sl2)}]
For $0<v<1$, $U_v(\mathfrak{sl}_2)$ is generated by
$E,F,K,K^{-1}$ with
\[
KEK^{-1}=v^2E,
\qquad KFK^{-1}=v^{-2}F,
\qquad [E,F]=\frac{K-K^{-1}}{v-v^{-1}},
\]
and
\[
\Delta(K)=K\otimes K,
\quad
\Delta(E)=E\otimes1+K\otimes E,
\quad
\Delta(F)=F\otimes K^{-1}+1\otimes F.
\]
Its counit, antipode, and compact real-form involution are
\[
\varepsilon(E)=\varepsilon(F)=0,
\quad \varepsilon(K)=1,
\quad
S(E)=-K^{-1}E,
\quad S(F)=-FK,
\quad S(K)=K^{-1},
\]
\[
K^*=K,
\qquad E^*=KF,
\qquad F^*=EK^{-1}.
\]
\end{definition}

\begin{exercise}[The \texorpdfstring{$U_v(\mathfrak{sl}_2)$}{Uv(sl2)} Hopf deformation]
Verify that the coproduct preserves the defining relations and that the
displayed counit, antipode, and involution make
$U_v(\mathfrak{sl}_2)$ a Hopf $*$-algebra.  In the $\hbar$-adic
completion put
\[
v=e^{-\hbar},\qquad K=v^H=e^{-\hbar H}.
\]
Prove that
\[
[H,E]=2E,\qquad [H,F]=-2F,
\qquad
\frac{K-K^{-1}}{v-v^{-1}}=H+O(\hbar^2),
\]
and that the relations and coproduct tend to those of
$U(\mathfrak{sl}_2)$ as $\hbar\to0$.  Expand the coproduct through first
order and compute the resulting coboundary Lie bialgebra.
\end{exercise}

\begin{definition}[Gaussian binomial coefficients]
For an indeterminate $u$, $n\geq0$, and $0\leq r\leq n$, define
\[
\genfrac{[}{]}{0pt}{}{n}{r}_{u}
=\prod_{j=1}^r\frac{1-u^{n-r+j}}{1-u^j},
\]
with value $1$ for $r=0$ and value $0$ outside this range.
\end{definition}

\begin{exercise}[Gaussian and quantum-group binomial identities]
Derive
\[
\genfrac{[}{]}{0pt}{}{n}{r}_{u}
=\genfrac{[}{]}{0pt}{}{n-1}{r}_{u}
+u^{n-r}\genfrac{[}{]}{0pt}{}{n-1}{r-1}_{u}
\]
and prove
\[
\prod_{j=0}^{n-1}(1+tu^j)
=\sum_{r=0}^n
u^{r(r-1)/2}\genfrac{[}{]}{0pt}{}{n}{r}_{u}t^r.
\]
For $YX=uXY$, prove the noncommutative identity
\[
(X+Y)^n=\sum_{r=0}^n
\genfrac{[}{]}{0pt}{}{n}{r}_{u}X^rY^{n-r}.
\]
Set $u=v^2$ and apply it to
$X=E\otimes1$ and $Y=K\otimes E$ in $U_v(\mathfrak{sl}_2)$ to obtain an
explicit formula for $\Delta(E^n)$.  Show that the limit $v\to1$ recovers
the ordinary binomial theorem and the undeformed coproduct.
\end{exercise}


\begin{exercise}[Gaussian-binomial state counts]
Apply the Gaussian binomial theorem from the preceding exercise.
Prove that the coefficient of $u^m$ in
$\genfrac{[}{]}{0pt}{}{n}{r}_u$ counts partitions
of $m$ whose diagrams fit inside an $r\times(n-r)$ rectangle.  For $n$
fermionic levels with energies $j\hbar\omega$, $0\leq j<n$, prove that the
coefficient of $t^ru^{m+r(r-1)/2}$ in
\[
\prod_{j=0}^{n-1}(1+tu^j)
\]
is the exact number of $r$-particle states at excitation energy
$m\hbar\omega$ above the filled lowest configuration, and identify it with
the coefficient of $u^m$ in $\genfrac{[}{]}{0pt}{}{n}{r}_u$.
\end{exercise}

\begin{definition}[Compact quantum groups and Haar states]
A \emph{compact quantum group} $\mathbb G$ is a unital $C^*$-algebra
$C(\mathbb G)$ with a unital $*$-homomorphism
\[
\Delta:C(\mathbb G)\longrightarrow
C(\mathbb G)\otimes_{\min}C(\mathbb G)
\]
where $\otimes_{\min}$ is the closure of the algebraic tensor product in
$B(H\otimes K)$ for faithful representations on $H$ and $K$.  The
coproduct satisfies
\[
(\Delta\otimes\id)\Delta=(\id\otimes\Delta)\Delta
\]
and
\[
\overline{\Delta(C(\mathbb G))(1\otimes C(\mathbb G))}
=C(\mathbb G)\otimes_{\min}C(\mathbb G)
=\overline{\Delta(C(\mathbb G))(C(\mathbb G)\otimes1)}.
\]
A \emph{Haar state} is a state
$h$ satisfying
\[
(h\otimes\id)\Delta(a)=h(a)1=(\id\otimes h)\Delta(a).
\]
\end{definition}

\begin{definition}[Representative functions]
A finite-dimensional \emph{unitary corepresentation} on $H_U$ is a unitary
$U\in B(H_U)\otimes C(\mathbb G)$ satisfying
\[
(\id\otimes\Delta)(U)=U_{12}U_{13}.
\]
For $\omega\in B(H_U)^*$, the corresponding \emph{coefficient} is
\[
u_\omega=(\omega\otimes\id)(U)\in C(\mathbb G).
\]
The representative-function algebra $\operatorname{Pol}(\mathbb G)$ is
the unital $*$-subalgebra generated by these coefficients for all
finite-dimensional unitary corepresentations.
\end{definition}

\begin{definition}[$SU_v(2)$]
Let $\operatorname{Pol}(SU_v(2))$ be the universal unital $*$-algebra
generated by $\alpha,\gamma$ so that
\[
u=\begin{pmatrix}
\alpha&-v\gamma^*\\
\gamma&\alpha^*
\end{pmatrix}
\]
is unitary.  Equivalently, its relations are
\[
\begin{gathered}
\alpha^*\alpha+\gamma^*\gamma=1,
\qquad \alpha\alpha^*+v^2\gamma\gamma^*=1,
\qquad \gamma^*\gamma=\gamma\gamma^*,\\
\alpha\gamma=v\gamma\alpha,
\qquad \alpha\gamma^*=v\gamma^*\alpha,
\end{gathered}
\]
Define $\Delta(u_{ij})=\sum_k u_{ik}\otimes u_{kj}$, and let
$C(SU_v(2))$ be the universal $C^*$-completion.
\end{definition}

\begin{exercise}[$SU_v(2)$ as a compact quantum group]
Prove that $\Delta$ is a coassociative $*$-homomorphism satisfying the
cancellation conditions.  On $\ell^2(\Z_{\geq0})$, define for $\xi\in\T$
\[
\pi_\xi(\gamma)e_n=\xi v^ne_n,
\qquad
\pi_\xi(\alpha)e_n=(1-v^{2n})^{1/2}e_{n-1},
\qquad e_{-1}=0.
\]
Verify the relations and prove that
\[
h(a)=(1-v^2)\sum_{n\geq0}v^{2n}
\int_\T\langle e_n,\pi_\xi(a)e_n\rangle\,dm_\T(\xi)
\]
is the Haar state.
\end{exercise}

\begin{definition}[Quantum-group coefficient pairing]
For a finite-dimensional type-$1$ $U_v(\mathfrak{sl}_2)$-module $V$,
$\xi\in V$, and $\varphi\in V^\vee$, define
\[
c_{\varphi,\xi}(x)=\varphi(x\xi)
\qquad(x\in U_v(\mathfrak{sl}_2)).
\]
Let $\mathcal O_v$ be the Hopf algebra generated by these coefficient
functionals.  Its pairing with $U_v(\mathfrak{sl}_2)$ is
\[
\langle x,c_{\varphi,\xi}\rangle=c_{\varphi,\xi}(x),
\]
with multiplication and coproduct characterized by
\[
\langle x,ab\rangle=\langle\Delta x,a\otimes b\rangle,
\qquad
\langle xy,a\rangle=\langle x\otimes y,\Delta a\rangle.
\]
\end{definition}

\begin{exercise}[Coefficient corepresentations]
Prove that the coefficient map of a tensor product is the product of the
coefficient maps, that the coefficient map of a dual module is obtained
from the antipode, and that the displayed pairing is a nondegenerate Hopf
pairing.  Apply it to the fundamental module and prove that sending its
coefficient corepresentation to the displayed matrix $u$ extends to a
Hopf $*$-isomorphism
\[
\mathcal O_v\longrightarrow\operatorname{Pol}(SU_v(2)).
\]
\end{exercise}

\begin{exercise}[$SU_v(2)$ Peter--Weyl decomposition]
Construct the irreducible type-$1$ unitary
$U_v(\mathfrak{sl}_2)$-modules $V_n$ from their highest-weight lines and
form their coefficient corepresentations of $\operatorname{Pol}(SU_v(2))$
using the Hopf pairing.
Use the explicit Haar state to derive quantum Schur orthogonality and prove
\[
\operatorname{Pol}(SU_v(2))
=\bigoplus_{n\geq0}C(V_n)
\]
algebraically and as a Hilbert direct sum in the Haar GNS space.
\end{exercise}

\begin{definition}[Conjugates and quantum dimension]
For a finite-dimensional $H$-module $V$, its \emph{left dual} $V^\vee$ has
action
\[
(x\varphi)(v)=\varphi(S(x)v),
\]
evaluation $\operatorname{ev}_V(\varphi\otimes v)=\varphi(v)$, and
coevaluation
$\operatorname{coev}_V$.  In a rigid
$C^*$-tensor category, a \emph{conjugate} of $V$ is an object
$\overline V$ with morphisms
\[
R_V:\mathbf 1\longrightarrow\overline V\otimes V,
\qquad
\overline R_V:\mathbf 1\longrightarrow V\otimes\overline V
\]
satisfying
\[
(\overline R_V^*\otimes\id_V)(\id_V\otimes R_V)=\id_V,
\qquad
(R_V^*\otimes\id_{\overline V})
(\id_{\overline V}\otimes\overline R_V)=\id_{\overline V}.
\]
The adjoints $R_V^*$ and $\overline R_V^*$ are evaluation maps, and
$R_V$ and $\overline R_V$ are coevaluation maps.  A solution is
\emph{standard} when
\[
R_V^*(\id_{\overline V}\otimes T)R_V
=\overline R_V^*(T\otimes\id_{\overline V})\overline R_V
\qquad(T\in\End(V)).
\]
For a standard solution, define
\[
\operatorname{tr}_v(T)
=R_V^*(\id_{\overline V}\otimes T)R_V,
\qquad
\dim_v(V)=\operatorname{tr}_v(\id_V).
\]
\end{definition}

\begin{exercise}[$SU_v(2)$ fusion and dimensions]
Let $V_n$ be the irreducible type-$1$ $U_v(\mathfrak{sl}_2)$-module of
highest weight $n$.  Construct its conjugate equations and prove
\[
V_m\otimes V_n
\cong\bigoplus_{k=0}^{\min(m,n)}V_{m+n-2k}.
\]
Normalize the conjugate solutions by the compact involution and compute
\[
\dim_v(V_n)=[n+1]_v
=\frac{v^{n+1}-v^{-(n+1)}}{v-v^{-1}}.
\]
Prove
\[
\dim V_n=n+1,
\qquad
\lim_{v\to1}\dim_v(V_n)=n+1,
\qquad
\dim_v(V_n)>n+1
\]
for $0<v<1$ and $n\geq1$.
\end{exercise}

\begin{definition}[Basic hypergeometric series]
For $m\in\N$, define
\[
(a;Q)_m=\prod_{j=0}^{m-1}(1-aQ^j),
\qquad
(a;Q)_\infty=\prod_{j=0}^{\infty}(1-aQ^j),
\]
\[
(a_1,\ldots,a_r;Q)_m=\prod_{i=1}^r(a_i;Q)_m,
\qquad
(a_1,\ldots,a_r;Q)_\infty=\prod_{i=1}^r(a_i;Q)_\infty.
\]
The \emph{basic hypergeometric series} is
\[
{}_r\phi_s\!\left(
\begin{matrix}a_1,\ldots,a_r\\b_1,\ldots,b_s\end{matrix};Q,z
\right)
=\sum_{m=0}^\infty
\frac{(a_1,\ldots,a_r;Q)_m}
     {(Q,b_1,\ldots,b_s;Q)_m}
\left((-1)^mQ^{m(m-1)/2}\right)^{1+s-r}z^m.
\]
It is terminating when one numerator parameter is $Q^{-n}$.
\end{definition}

\begin{exercise}[Basic hypergeometric summations]
Prove the $Q$-Gauss identity
\[
{}_2\phi_1\!\left(
\begin{matrix}a,b\\c\end{matrix};Q,\frac{c}{ab}
\right)
=\frac{(c/a,c/b;Q)_\infty}{(c,c/(ab);Q)_\infty},
\qquad \left|\frac{c}{ab}\right|<1,
\]
and the terminating $Q$-Pfaff--Saalsch\"utz identity
\[
{}_3\phi_2\!\left(
\begin{matrix}a,b,Q^{-n}\\c,abQ^{1-n}/c\end{matrix};Q,Q
\right)
=\frac{(c/a,c/b;Q)_n}{(c,c/(ab);Q)_n}
\]
by comparing coefficients after multiplication by the denominator
$Q$-products in the first identity, and by showing that the two sides of
the second identity satisfy the same recurrence in $n$ and have the same
initial value.
\end{exercise}

\begin{definition}[Quantum coupling coefficients]
For irreducible type-$1$ modules $V_a,V_b,V_c$, let
\[
C_{ab}^c:V_c\longrightarrow V_a\otimes V_b
\]
be the intertwining isometry normalized by its positive highest-weight
coefficient.  Its matrix coefficients between the distinguished weight
lines are the \emph{$v$-Clebsch--Gordan coefficients}.
\end{definition}

\begin{exercise}[XXZ quantum-group symmetry]
Use the distinguished weight-line realization
$V_1=\C e_+\oplus\C e_-$ of the fundamental type-$1$ module.  Define the
Pauli endomorphisms, relative to the ordered weight lines, by
\[
\sigma^x=\begin{pmatrix}0&1\\1&0\end{pmatrix},
\qquad
\sigma^y=\begin{pmatrix}0&-i\\i&0\end{pmatrix},
\qquad
\sigma^z=\begin{pmatrix}1&0\\0&-1\end{pmatrix},
\]
and let $\sigma_j^a$ act as $\sigma^a$ on the $j$th tensor factor and as
the identity on the others.  On $V_1^{\otimes L}$ let
\[
\begin{aligned}
H_{\rm XXZ}=J\sum_{j=1}^{L-1}\biggl(&
\sigma_j^x\sigma_{j+1}^x+\sigma_j^y\sigma_{j+1}^y
+\frac{v+v^{-1}}2(\sigma_j^z\sigma_{j+1}^z-I)\\
&+\frac{v-v^{-1}}2(\sigma_j^z-\sigma_{j+1}^z)\biggr).
\end{aligned}
\]
For $J\in\R$ and $0<v<1$, prove that this open-chain Hamiltonian is
self-adjoint and commutes with the iterated coproduct action of
$U_v(\mathfrak{sl}_2)$.
\end{exercise}

\begin{exercise}[Dual \texorpdfstring{$Q$}{Q}-Hahn coupling coefficients]
Decompose two coupled irreducible spin sectors by
the projections $C_{ab}^c(C_{ab}^c)^*$ and derive the orthogonality
relations for their $v$-Clebsch--Gordan coefficients.  Prove that their
weight dependence satisfies the dual $Q$-Hahn recurrence with $Q=v^2$ and
normalized amplitudes with a terminating ${}_3\phi_2$ of base $Q$.

\end{exercise}

\begin{exercise}[XXZ transition strengths]
Let the weak drive be an irreducible tensor operator $A^{(s)}_\mu$ for the
$U_v(\mathfrak{sl}_2)$ action.  Prove the $v$-Wigner--Eckart factorization
of its matrix elements into a reduced matrix element and a
$v$-Clebsch--Gordan coefficient.  Derive the relative transition strengths
at fixed reduced matrix element as squared $v$-Clebsch--Gordan
coefficients.  State
the chain length, open boundary terms, resolved weights, energy scale $J$,
and nondegeneracy assumptions under which the measured line positions and
strengths determine $v$.
\end{exercise}

\begin{definition}[Universal $R$-matrix and braiding]
A \emph{universal $R$-matrix} for a Hopf algebra $H$ is an invertible
element $\mathcal R$ in a completion of $H\otimes H$ on which the
finite-dimensional tensor-product representations under consideration
extend, satisfying
\[
\mathcal R\Delta(x)\mathcal R^{-1}=\Delta^{\rm op}(x),
\quad
(\Delta\otimes\id)(\mathcal R)=\mathcal R_{13}\mathcal R_{23},
\quad
(\id\otimes\Delta)(\mathcal R)=\mathcal R_{13}\mathcal R_{12}.
\]
For finite-dimensional modules $V,W$, define
\[
c_{V,W}
=\tau\circ(\pi_V\otimes\pi_W)(\mathcal R):
V\otimes W\longrightarrow W\otimes V.
\]
The pair $(H,\mathcal R)$ is \emph{quasitriangular}.
\end{definition}

\begin{definition}[Spectral exchange operators]
Let $V$ be finite-dimensional.  A meromorphic family
$R(u)\in\End(V\otimes V)$ satisfies the spectral Yang--Baxter equation
when
\[
R_{12}(u/v)R_{13}(u/w)R_{23}(v/w)
=R_{23}(v/w)R_{13}(u/w)R_{12}(u/v).
\]
It is regular when $R(1)$ is a nonzero scalar multiple of the flip and
unitary up to scalar when $R_{12}(u)R_{21}(u^{-1})$ is scalar.
\end{definition}

\begin{exercise}[Hexagon identities]
Prove that $c_{V,W}$ is a natural module intertwiner and that the two
coproduct identities for $\mathcal R$ give the two hexagon identities.
Derive
\[
\mathcal R_{12}\mathcal R_{13}\mathcal R_{23}
=\mathcal R_{23}\mathcal R_{13}\mathcal R_{12}
\]
and, on $V^{\otimes3}$,
\[
c_{12}c_{23}c_{12}=c_{23}c_{12}c_{23}.
\]
Prove that the operators $c_{j,j+1}$ define braid-group representations on
tensor powers of every finite-dimensional type-$1$ module.
\end{exercise}

\begin{definition}[Ribbon categories]
Let $\mathcal C$ be a braided rigid $\C$-linear monoidal category.  A
\emph{twist} is a natural isomorphism
$\theta:\id_{\mathcal C}\to\id_{\mathcal C}$ satisfying
\[
\theta_{\mathbf1}=\id_{\mathbf1},
\qquad
\theta_{X\otimes Y}
=c_{Y,X}c_{X,Y}(\theta_X\otimes\theta_Y).
\]
The category is \emph{ribbon} when
\[
\theta_{X^\vee}=(\theta_X)^\vee
\]
for every object $X$.
\end{definition}

\begin{exercise}[Ribbon elements]
For a quasitriangular Hopf algebra $(H,\mathcal R)$, set
\[
u=m(S\otimes\id)(\mathcal R_{21}).
\]
Let $\nu_R\in H$ be central and invertible and satisfy
\[
\nu_R^2=uS(u),
\qquad S(\nu_R)=\nu_R,
\qquad \varepsilon(\nu_R)=1,
\qquad
\Delta(\nu_R)=(\mathcal R_{21}\mathcal R)^{-1}
(\nu_R\otimes \nu_R).
\]
Prove that
\[
\theta_V=\pi_V(\nu_R^{-1})
\]
defines a ribbon twist on the category of finite-dimensional $H$-modules.
Derive invariance of the quantum trace under framed ribbon isotopy.
\end{exercise}

\begin{exercise}[Chebyshev quantum dimensions]
For the irreducible $SU_v(2)$ objects $V_n$, use
$V_1\otimes V_n\cong V_{n+1}\oplus V_{n-1}$ to prove
\[
d_{n+1}=(v+v^{-1})d_n-d_{n-1},
\qquad d_0=1,\quad d_1=v+v^{-1},
\]
where $d_n=\dim_v(V_n)$.  If $U_n$ is the Chebyshev polynomial determined
by $U_0(x)=1$, $U_1(x)=2x$, and
$U_{n+1}(x)=2xU_n(x)-U_{n-1}(x)$, deduce
\[
\dim_v(V_n)=U_n\!\left(\frac{v+v^{-1}}2\right)
=\frac{v^{n+1}-v^{-(n+1)}}{v-v^{-1}}.
\]
\end{exercise}

\begin{definition}[Quantum recoupling coefficients]
The matrix coefficients of the recoupling map
\[
(V_a\otimes V_b)\otimes V_c
\longrightarrow
V_a\otimes(V_b\otimes V_c)
\]
between fixed intermediate irreducible summands are the $v$-$6j$
coefficients.
\end{definition}

\begin{exercise}[Quantum Racah recoupling]
Compute the change between the two coupling orders of
$V_a\otimes V_b\otimes V_c$ and express its entries as terminating
${}_4\phi_3$ series with base $Q=v^2$.  Identify them with
$Q$-Racah polynomials and derive
their orthogonality from unitarity of the recoupling map for the compact
real form.  Prove the pentagon identity by comparing the five coupling
orders of four factors.  For three resolved XXZ spin sectors, derive the
recoupling probabilities as squared $v$-$6j$ coefficients and state the
unitarity and spectral-resolution assumptions required for this
probability interpretation.
\end{exercise}



\begin{definition}[Type-$A$ Hecke actions]
For $v\ne0$, the Hecke algebra $\mathcal H_v(S_r)$ is generated by
$T_1,\ldots,T_{r-1}$ with
\[
T_iT_{i+1}T_i=T_{i+1}T_iT_{i+1},
\qquad T_iT_j=T_jT_i\quad(|i-j|>1),
\]
\[
(T_i-v)(T_i+v^{-1})=0.
\]
Let $V_1$ be the fundamental $U_v(\mathfrak{sl}_2)$-module.  On
$V_1\otimes V_1\cong V_2\oplus V_0$, let $P_2,P_0$ be the invariant
orthogonal projections, and define
\[
T=vP_2-v^{-1}P_0,
\qquad
\rho_r(T_i)=\id^{\otimes(i-1)}\otimes T
 \otimes\id^{\otimes(r-i-1)}.
\]
\end{definition}

\begin{exercise}[Fundamental Hecke symmetry]
Prove that the coproduct action of $U_v(\mathfrak{sl}_2)$ commutes with
$\rho_r(T_i)$.  Derive the braid relation from the $Q$-Racah recoupling
identity and the Hecke relation from the two displayed spectral
projections.  Put $e_i=(v-T_i)/(v+v^{-1})$ and prove that the action factors
through the Temperley--Lieb relations
\[
e_i^2=e_i,
\qquad e_ie_{i\pm1}e_i=(v+v^{-1})^{-2}e_i.
\]
\end{exercise}

\begin{exercise}[Hecke Poincar\'e identity]
Put $\tau=v^2$.  For $w\in S_r$, let $\ell(w)$ be the least number of adjacent
transpositions in an expression for $w$, and let
$T_w=T_{i_1}\cdots T_{i_{\ell(w)}}$ for any reduced expression
$w=s_{i_1}\cdots s_{i_{\ell(w)}}$.  Prove from the braid relations that
$T_w$ is well-defined and that
\[
\sum_{w\in S_r}\tau^{\ell(w)}
=\prod_{j=1}^r\frac{1-\tau^j}{1-\tau}.
\]
Prove that
\[
e_+=\frac{\sum_{w\in S_r}v^{\ell(w)}T_w}
{\sum_{w\in S_r}v^{2\ell(w)}},
\qquad
e_-=\frac{\sum_{w\in S_r}(-v^{-1})^{\ell(w)}T_w}
{\sum_{w\in S_r}v^{-2\ell(w)}}
\]
are idempotents satisfying
$T_ie_+=ve_+$ and $T_ie_-=-v^{-1}e_-$.  Under the fundamental Hecke action
from the preceding exercise, identify their images with the quantum symmetric
and exterior state spaces and interpret the Poincar\'e polynomial as their
graded state count.
\end{exercise}

\begin{definition}[Affine Hecke algebras]
The extended affine Hecke algebra
$\mathcal H_v^{\mathrm{aff}}(GL_r)$ is generated by
$T_1,\ldots,T_{r-1}$ and commuting invertible elements
$X_1,\ldots,X_r$, subject to the finite Hecke relations and
\[
T_iX_iT_i=X_{i+1},
\qquad
T_iX_j=X_jT_i\quad(j\neq i,i+1).
\]
Equivalently, it is the Hecke quotient of the extended affine braid group
of $\Z^r\rtimes S_r$.
\end{definition}

\begin{definition}[Baxterized Hecke operators]
For an invertible spectral parameter $u$, define
\[
\check R_i(u)=uT_i-u^{-1}T_i^{-1}.
\]
Its eigenvalues on the $v$ and $-v^{-1}$ eigenspaces of $T_i$ are
\[
uv-u^{-1}v^{-1},
\qquad
-uv^{-1}+u^{-1}v.
\]
\end{definition}

\begin{exercise}[Baxterization]
Use the braid and Hecke relations for $T_i$ to prove
\[
\check R_i(u)\check R_{i+1}(uv)\check R_i(v)
=\check R_{i+1}(v)\check R_i(uv)\check R_{i+1}(u).
\]
Prove that $\check R_i(1)$ is a scalar multiple of the identity and compute
$\check R_i(u)\check R_i(u^{-1})$.  For $|u|=1$, prove that
$\check R_i(u)^*=\check R_i(u^{-1})$ and normalize the operators to be
unitary wherever the scalar product is nonzero.  Prove that their two
eigenvalues depend on $u$ and hence that the spectral family itself does
not satisfy one parameter-independent Hecke polynomial.
\end{exercise}

\begin{definition}[$q$KZ transport]
Let $v$ be the quantum-group parameter and let $\kappa\ne0$ be an
independent shift parameter.  Let $D\in\operatorname{End}(V)$ be
invertible and satisfy $[D\otimes D,\check R(u)]=0$.  Let $D_1$ act on
the transported tensor factor and define
\[
\mathcal K_1(z_1,\ldots,z_r)
=D_1\check R_{1r}(z_1/z_r)\cdots\check R_{12}(z_1/z_2).
\]
The first $q$KZ equation is
\[
\Psi(z_2,\ldots,z_r,\kappa z_1)
=\mathcal K_1(z_1,\ldots,z_r)\Psi(z_1,\ldots,z_r),
\]
and the remaining equations are obtained by cyclic conjugation.
\end{definition}

\begin{exercise}[$q$KZ exchange transport]
Use Baxterization to prove consistency of the adjacent exchange rules and
the $q$KZ difference equations.  Identify the constant exchanges and
winding operators with the generators of
$\mathcal H_v^{\mathrm{aff}}(GL_r)$, keeping $v$ and $\kappa$ independent.
For the unitary normalization, prepare and detect three resolved XXZ spin
sectors and prove that the two factorizations of a three-spin exchange give
the same output probabilities.  Express this equality as the $Q$-Racah
recoupling identity.  Derive the phase accumulated in one $\kappa$-shift
cycle and state the twist, unitarity, preparation, and detection assumptions
under which it is an interference observable.
\end{exercise}

\begin{definition}[Temperley--Lieb algebras]
For $d\in\C$, the Temperley--Lieb algebra $TL_n(d)$ is generated by
$U_1,\ldots,U_{n-1}$ with
\[
U_i^2=dU_i,
\qquad
U_iU_{i\pm1}U_i=U_i,
\qquad
U_iU_j=U_jU_i\quad(|i-j|>1).
\]
Its diagram trace is normalized by
\[
\operatorname{tr}_0(1)=1,
\qquad
\operatorname{tr}_{n+1}(x)=d\operatorname{tr}_n(x),
\qquad
\operatorname{tr}_{n+1}(xU_n)=\operatorname{tr}_n(x),
\]
where $x\in TL_n(d)$ is embedded in $TL_{n+1}(d)$.
\end{definition}

\begin{definition}[Jones--Wenzl recursion]
Put
\[
\zeta=e^{\pi i/(k+2)},
\qquad
d=\zeta+\zeta^{-1},
\qquad
[n]_\zeta=\frac{\zeta^n-\zeta^{-n}}{\zeta-\zeta^{-1}}.
\]
With $f_0=f_1=1$, define recursively in $TL_{n+1}(d)$
\[
f_{n+1}=f_n-\frac{[n]_\zeta}{[n+1]_\zeta}f_nU_nf_n
\qquad(1\leq n\leq k).
\]
\end{definition}

\begin{definition}[Negligible morphisms]
In a pivotal category with trace $\operatorname{tr}$, a morphism
$f:X\to Y$ is \emph{negligible} when
\[
\operatorname{tr}(gf)=0
\qquad\text{for every }g:Y\to X.
\]
Let $\mathcal N$ be the tensor ideal of negligible morphisms.  Let
$\mathcal C_k$ be the idempotent completion of the quotient of the
Temperley--Lieb category at $d=\zeta+\zeta^{-1}$ by $\mathcal N$, and write
$V_a$ for the image of $f_a$, $0\leq a\leq k$.
\end{definition}

\begin{exercise}[Jones--Wenzl trace]
Use the Temperley--Lieb relations to prove inductively
\[
f_n^2=f_n,
\qquad
U_if_n=f_nU_i=0\quad(1\leq i<n),
\qquad
\operatorname{tr}(f_n)=[n+1]_\zeta.
\]
Deduce $\operatorname{tr}(f_{k+1})=[k+2]_\zeta=0$.
\end{exercise}

\begin{exercise}[Temperley--Lieb semisimplification]
Prove that
$\mathcal N$ is the tensor ideal generated by $f_{k+1}$, and prove that
the quotient removes every trace-zero summand and is finite semisimple.
\end{exercise}

\begin{exercise}[Truncated fusion]
Derive
\[
V_a\otimes V_b
\cong
\bigoplus_{\substack{|a-b|\leq c\leq
\min(a+b,2k-a-b)\\a+b+c\equiv0\pmod2}}V_c
\]
and
\[
d_a=[a+1]_\zeta
=\frac{\sin((a+1)\pi/(k+2))}{\sin(\pi/(k+2))}.
\]
\end{exercise}

\begin{definition}[\texorpdfstring{$SU(2)_k$}{SU(2)k} modular data]
For $0\leq a,b\leq k$, define
\[
S_{ab}
=\sqrt{\frac{2}{k+2}}
\sin\!\left(\frac{(a+1)(b+1)\pi}{k+2}\right),
\qquad
\theta_a
=\exp\!\left(\frac{\pi i\,a(a+2)}{2(k+2)}\right).
\]
Let $N_{ab}^{\ \ c}$ be the multiplicity of $V_c$ in
$V_a\otimes V_b$.
\end{definition}

\begin{exercise}[Verlinde diagonalization]
Use sine orthogonality to prove that $S$ is real, symmetric, and unitary.
For each $a$, let $(N_a)_{bc}=N_{ab}^{\ \ c}$.  Use the truncated
Chebyshev recursion to prove
\[
N_a=S\,\operatorname{diag}\!\left(
\frac{S_{ax}}{S_{0x}}
\right)_{x=0}^kS^{-1}
\]
and hence
\[
N_{ab}^{\ \ c}
=\sum_{x=0}^k
\frac{S_{ax}S_{bx}\overline{S_{cx}}}{S_{0x}}.
\]
Deduce that the number of fusion channels from
$a_1,\ldots,a_n$ to $c$ is
$(N_{a_1}\cdots N_{a_n})_{0c}$ and express it as a finite sum of products
of sine functions.  Apply the same formula to the Ising modular matrix
derived in Counting and recover the Ising fusion multiplicities.
\end{exercise}

\begin{exercise}[Anyon interference and fusion counting]
For simple charges $a,b$, prove that the normalized monodromy amplitude is
\[
M_{ab}=\frac{S_{ab}S_{00}}{S_{0a}S_{0b}}.
\]
In a two-path anyon interferometer with path intensities $I_1,I_2$, derive
\[
I_b(\phi)=I_1+I_2
+2\sqrt{I_1I_2}\,\Re(e^{i\phi}M_{ab}).
\]
Insert the sine formula and predict the visibility and phase for every
enclosed charge.  For a collection of prepared charges, use the Verlinde
formula to predict the exact number of fusion outcomes in each total-charge
sector.  Identify indistinguishable charges for one probe and prove that
the full family of probes distinguishes them.  State the coherence,
topological-gap, charge-preparation, and readout assumptions separately
from the predicted fringes and integer state counts.
\end{exercise}

\section{Integrable Systems}

\begin{exercise}[Elliptic pendulum period]
On $T^*S^1$ consider
\[
H(\theta,p)=\frac{p^2}{2m\ell^2}+mg\ell(1-\cos\theta).
\]
For an oscillation of amplitude $0<\theta_0<\pi$, set
$k=\sin(\theta_0/2)$ and define
\[
K(k)=\int_0^{\pi/2}\frac{d\phi}
{\sqrt{1-k^2\sin^2\phi}}.
\]
Use energy conservation to prove
\[
T(\theta_0)=4\sqrt{\frac{\ell}{g}}\,K(k).
\]
Apply the generalized binomial theorem with $a=\frac12$.  For
\[
I_m=\int_0^{\pi/2}\sin^{2m}\phi\,d\phi,
\]
use integration by parts to prove
\[
I_0=\frac{\pi}{2},
\qquad
I_m=\frac{2m-1}{2m}I_{m-1}\quad(m\geq1),
\]
and obtain
\[
K(k)=\frac{\pi}{2}\sum_{m=0}^{\infty}
\left(\frac{(1/2)_m}{m!}\right)^2k^{2m},
\qquad |k|<1.
\]
Deduce
\[
T(\theta_0)=2\pi\sqrt{\frac{\ell}{g}}
\left(1+\frac{\theta_0^2}{16}+O(\theta_0^4)\right).
\]
Deduce the experimentally testable increase of period with amplitude and
the logarithmic divergence as $\theta_0\uparrow\pi$.
\end{exercise}

\begin{definition}[Sutherland model]
For $N$ identical particles of mass $m$ on a ring of radius $R$ and
$g\geq1$, let
\[
\operatorname{Conf}_N(S^1)
=\{(\theta_1,\ldots,\theta_N):\theta_i\ne\theta_j\text{ for }i\ne j\}.
\]
On $C_c^\infty(\operatorname{Conf}_N(S^1))^{S_N}$ define
\[
\widehat H_g=\frac{\hbar^2}{2mR^2}
\left[-\sum_{j=1}^N\partial_{\theta_j}^2
+\frac{g(g-1)}2\sum_{i<j}
\frac1{\sin^2((\theta_i-\theta_j)/2)}\right]
\]
and let the same symbol denote its Friedrichs realization in
$L^2((S^1)^N)^{S_N}$.  Define
\[
\Psi_0(\theta)=\prod_{i<j}|e^{i\theta_i}-e^{i\theta_j}|^g.
\]
For a partition $\lambda$ of length at most $N$, use the monomial
symmetric polynomial $m_\lambda$ and dominance order defined in Counting,
evaluated at $z_j=e^{i\theta_j}$.
\end{definition}

\begin{definition}[Jack polynomials]
The Jack polynomial $P_\lambda^{(1/g)}$ is the symmetric polynomial
normalized by
\[
P_\lambda^{(1/g)}=m_\lambda+
\sum_{\mu<\lambda}c_{\lambda\mu}(g)m_\mu
\]
and by satisfying
\[
\left[
\sum_i(z_i\partial_{z_i})^2
+g\sum_{i<j}\frac{z_i+z_j}{z_i-z_j}
(z_i\partial_{z_i}-z_j\partial_{z_j})
\right]P_\lambda^{(1/g)}
=\left[\sum_i
\bigl(\lambda_i^2+g(N+1-2i)\lambda_i\bigr)\right]
P_\lambda^{(1/g)};
\]
here $<$ is dominance order.
\end{definition}

\begin{exercise}[Sutherland spectrum]
Conjugate the quantum Hamiltonian by $\Psi_0$, identify its commuting
quantum integrals, and prove that
\[
\Psi_\lambda=\Psi_0P_\lambda^{(1/g)}(e^{i\theta_1},\ldots,e^{i\theta_N})
\]
is a joint eigenfunction with
\[
E_\lambda=\frac{\hbar^2}{2mR^2}
\left[\frac{g^2N(N^2-1)}{12}
+\sum_{j=1}^N
\bigl(\lambda_j^2+g(N+1-2j)\lambda_j\bigr)\right].
\]
At $g=1$, prove that $P_\lambda^{(1)}=s_\lambda$ for the Schur polynomial
defined in Counting.
\end{exercise}

\begin{exercise}[Jack transition spectrum]
Let $\rho_r=\sum_{j=1}^Ne^{ir\theta_j}$ be a density mode.  Derive the
Jack Pieri identity in orthonormal normalization for
$\rho_1P_\lambda^{(1/g)}$ and show that only
nonincreasing tuples obtained from $\lambda$ by increasing one entry by
one occur.
For a weak periodic drive proportional to
$\rho_1e^{-i\Omega t}+\rho_{-1}e^{i\Omega t}$, deduce the allowed
resonance frequencies
\[
\Omega_{\mu\lambda}=\frac{|E_\mu-E_\lambda|}{\hbar}
\]
and express their relative intensities as squared normalized Jack Pieri
coefficients.  State how spectroscopy of particles confined to a ring
determines $g$ from the line shifts and tests the Jack selection rule from
the missing lines.
\end{exercise}

\begin{definition}[Double-affine braid group]
For $r\geq2$, the \emph{double-affine braid group} $\widetilde B_r$ of
type $GL_r$ is generated by $T_i$ for $1\leq i<r$, commuting invertible
elements $X_j$,
commuting invertible elements $Y_j$, and a central invertible element
$\mathsf q$, subject to the braid relations and
\[
T_iX_iT_i=X_{i+1},\qquad
T_iX_j=X_jT_i\quad(j\neq i,i+1),
\]
\[
T_i^{-1}Y_iT_i^{-1}=Y_{i+1},\qquad
T_iY_j=Y_jT_i\quad(j\neq i,i+1),
\]
\[
Y_1X_1\cdots X_r
=\mathsf q\,X_1\cdots X_rY_1,
\qquad
X_1^{-1}Y_2=Y_2X_1^{-1}T_1^{-2}.
\]
\end{definition}

\begin{exercise}[From affine to double-affine braids]
Let $T^2=\R^2/\Z^2$.  Interpret $T_i$ as an adjacent braid and $X_j,Y_j$
as motions of the $j$th point around the two cycles of $T^2$.  Derive the
displayed relations from these motions and identify $\widetilde B_r$ with
the central extension of
\[
\pi_1\bigl(\operatorname{Conf}_r(T^2)/S_r\bigr)
\]
whose total cycle commutator is $\mathsf q$.  Prove that deleting either
the $X$-family or the $Y$-family leaves an affine braid group.
\end{exercise}

\begin{definition}[Type-$A$ double affine Hecke algebra]
For nonzero parameters $q,t$, the \emph{double affine Hecke algebra}
$\mathsf H_{q,t}(GL_r)$ is the quotient of
$\C(q,t)[\widetilde B_r]$ by
\[
\mathsf q=q,
\qquad
(T_i-t^{1/2})(T_i+t^{-1/2})=0.
\]
\end{definition}

\begin{definition}[Rank-one DAHA]
For nonzero $k_0,k_1,u_0,u_1$ and a choice of $q^{1/2}$, the rank-one DAHA
$\mathsf H_q(C_1^\vee,C_1)$ is generated by
$T_0^{\pm1},T_1^{\pm1},(T_0^\vee)^{\pm1},(T_1^\vee)^{\pm1}$ subject to
\[
(T_i-k_i)(T_i+k_i^{-1})=0,
\qquad
(T_i^\vee-u_i)(T_i^\vee+u_i^{-1})=0
\quad(i=0,1),
\]
\[
\qquad
T_1^\vee T_1T_0T_0^\vee=q^{-1/2}.
\]
For the finite Hecke generator $T_1$, define
\[
e=\frac{T_1+k_1^{-1}}{k_1+k_1^{-1}}.
\]
\end{definition}

\begin{definition}[Rank-one polynomial representation]
Fix nonzero $k_0,k_1,u_0,u_1$ and choose $q^{1/2}$.  On
$\mathcal L=\C[z^{\pm1}]$ put
\[
(s_1f)(z)=f(z^{-1}),\qquad
(s_0f)(z)=f(qz^{-1}),
\]
\[
\begin{aligned}
(T_1f)(z)
&=k_1f(z^{-1})
+\frac{(k_1-k_1^{-1})+(u_1-u_1^{-1})z}
{1-z^2}\bigl(f(z)-f(z^{-1})\bigr),\\
(T_0f)(z)
&=k_0f(qz^{-1})
+\frac{(k_0-k_0^{-1})+(u_0-u_0^{-1})q^{1/2}z^{-1}}
{1-qz^{-2}}\bigl(f(z)-f(qz^{-1})\bigr).
\end{aligned}
\]
Let $X$ act by multiplication by $z$.  The Askey--Wilson parameters in
this convention are
\[
a=k_1u_1,\qquad b=-k_1u_1^{-1},\qquad
c=q^{1/2}k_0u_0,\qquad d=-q^{1/2}k_0u_0^{-1}.
\]
Put
\[
T_1^\vee=X^{-1}T_1^{-1},
\qquad
T_0^\vee=q^{-1/2}T_0^{-1}X.
\]
\end{definition}

\begin{definition}[Generalized hypergeometric series]
If no denominator parameter $b_j$ is a nonpositive integer, define
\[
{}_rF_s\!\left(
\begin{matrix}a_1,\ldots,a_r\\b_1,\ldots,b_s\end{matrix};z
\right)
=\sum_{m=0}^\infty
\frac{(a_1)_m\cdots(a_r)_m}
{(b_1)_m\cdots(b_s)_m}\frac{z^m}{m!}.
\]
The series is finite when a numerator parameter is a nonpositive integer.
Otherwise it converges for every $z$ when $r\leq s$, for $|z|<1$ when
$r=s+1$, and only at $z=0$ when $r>s+1$.
\end{definition}

\begin{definition}[Askey--Wilson operator and polynomials]
For $0<q<1$ and nonzero $a,b,c,d$, define
\[
A(z)=\frac{(1-az)(1-bz)(1-cz)(1-dz)}
{(1-z^2)(1-qz^2)}
\]
and, on $f(z)=f(z^{-1})$,
\[
(\mathcal D_{AW}f)(z)
=A(z)\bigl(f(qz)-f(z)\bigr)
+A(z^{-1})\bigl(f(q^{-1}z)-f(z)\bigr).
\]
For $x=(z+z^{-1})/2$, the \emph{Askey--Wilson polynomial} is
\[
\begin{aligned}
p_n(x;a,b,c,d\mid q)
={}&a^{-n}(ab,ac,ad;q)_n\\
&\times{}_4\phi_3\!\left(
\begin{matrix}q^{-n},abcdq^{n-1},az,az^{-1}\\
ab,ac,ad\end{matrix};q,q
\right).
\end{aligned}
\]
Its weight on the unit circle is
\[
w_{AW}(z)=
\frac{(z^2,z^{-2};q)_\infty}
{(az,a/z,bz,b/z,cz,c/z,dz,d/z;q)_\infty}.
\]
\end{definition}

\begin{exercise}[Rank-one spherical operator]
Verify that the displayed reflection--difference operators preserve
$\mathcal L$, that $T_0,T_1,T_0^\vee,T_1^\vee$ satisfy the four quadratic
Hecke relations and the product relation, and that they give the rank-one
DAHA polynomial representation with the stated parameter map.
Prove
\[
e\mathcal L=\mathcal L^{z\mapsto z^{-1}}
=\C[z+z^{-1}],
\]
and prove that the spherical action contains multiplication by
$z+z^{-1}$ and $\mathcal D_{AW}$.
\end{exercise}

\begin{exercise}[Askey--Wilson spectral theorem]
Prove
\[
\mathcal D_{AW}p_n
=(q^{-n}-1)(1-abcdq^{n-1})p_n.
\]
For $a,b,c,d\in(-1,1)$, prove that $\mathcal D_{AW}$ is symmetric for
\[
\langle f,g\rangle_{AW}
=\int_{|z|=1}f(z)\overline{g(z)}w_{AW}(z)
\frac{dz}{2\pi iz},
\]
and deduce the orthogonality and three-term recurrence of the
Askey--Wilson polynomials.
\end{exercise}

\begin{definition}[Wilson polynomials]
For positive $a,b,c,d$, define
\[
\begin{aligned}
W_n(x^2;a,b,c,d)
={}&(a+b)_n(a+c)_n(a+d)_n\\
&\times{}_4F_3\!\left(
\begin{matrix}-n,n+a+b+c+d-1,a+ix,a-ix\\
a+b,a+c,a+d\end{matrix};1
\right).
\end{aligned}
\]
\end{definition}

\begin{exercise}[Wilson degeneration]
Set $q=e^{-\varepsilon}$ and scale the Askey--Wilson parameters and
spectral variable by
\[
(a,b,c,d,z)=(q^A,q^B,q^C,q^D,q^{ix}).
\]
After the scalar normalization by $(1-q)^{-3n}$, prove that the
Askey--Wilson polynomial tends to
$W_n(x^2;A,B,C,D)$ as $\varepsilon\downarrow0$.  Derive the Wilson
difference equation, weight, orthogonality, and three-term recurrence as
limits of the corresponding Askey--Wilson identities.
\end{exercise}


\begin{definition}[Type-$A$ polynomial representation]
Let
\[
\mathcal P_r=\C(q,t^{1/2})[x_1^{\pm1},\ldots,x_r^{\pm1}],
\]
let $X_j$ act by multiplication by $x_j$, and let $s_i$ exchange
$x_i$ and $x_{i+1}$.  Define
\[
T_i=t^{1/2}s_i+(t^{1/2}-t^{-1/2})
       \frac{s_i-1}{x_i/x_{i+1}-1},
\qquad 1\leq i<r,
\]
and
\[
(\pi f)(x_1,\ldots,x_r)=f(x_2,\ldots,x_r,qx_1).
\]
For generic independent parameters $q,t$, call the operator algebra
generated by $X_j^{\pm1}$, $T_i$, and $\pi^{\pm1}$ the type-$A$
polynomial operator model.
\end{definition}

\begin{exercise}[DAHA polynomial quotient]
Prove that the displayed operators preserve $\mathcal P_r$ and satisfy
the braid and Hecke relations.  Prove
\[
\pi X_i\pi^{-1}=X_{i+1}\quad(1\leq i<r),
\qquad
\pi X_r\pi^{-1}=qX_1,
\]
and identify the two affine-Hecke subalgebras generated respectively by
$T_i,X_j$ and by $T_i,Y_j$.  Prove that this representation is obtained
from the double-affine braid-group algebra by imposing
\[
(T_i-t^{1/2})(T_i+t^{-1/2})=0
\]
and specializing its central parameter to $q$.
Conclude that the polynomial operator model is a realization of
$\mathsf H_{q,t}(GL_r)$.
\end{exercise}

\begin{definition}[Macdonald operators and polynomials]
Let
\[
(T_{q,x_i}f)(x_1,\ldots,x_r)
=f(x_1,\ldots,qx_i,\ldots,x_r).
\]
The first type-$A$ \emph{Macdonald operator} is
\[
D_{q,t}=\sum_{i=1}^r
 \left(\prod_{j\ne i}\frac{tx_i-x_j}{x_i-x_j}\right)T_{q,x_i}.
\]
For generic $q,t$, the \emph{Macdonald polynomial}
$P_\lambda(x;q,t)$ is the symmetric polynomial triangular in the
monomial symmetric functions and satisfying
\[
D_{q,t}P_\lambda
=\left(\sum_{i=1}^r q^{\lambda_i}t^{r-i}\right)P_\lambda,
\qquad
\lambda_1\geq\cdots\geq\lambda_r\geq0.
\]
For $0<q,t<1$, define
\[
\Delta_{q,t}(x)=
\prod_{i\neq j}\frac{(x_i/x_j;q)_\infty}
{(t x_i/x_j;q)_\infty},
\]
and
\[
\langle f,g\rangle_{q,t}
=\frac1{r!}\int_{(S^1)^r}
f(x)\overline{g(x)}\Delta_{q,t}(x)
\prod_{j=1}^r\frac{dx_j}{2\pi i x_j}.
\]
\end{definition}

\begin{exercise}[Cherednik operators]
Define
\[
Y_i=T_i\cdots T_{r-1}\pi
T_1^{-1}\cdots T_{i-1}^{-1}.
\]
Use the displayed polynomial representation and the affine braid relations
to prove that the $Y_i$ commute.  Prove that their symmetric Laurent
polynomials preserve $\mathcal P_r^{S_r}$ and identify
$t^{-(r-1)/2}(Y_1+\cdots+Y_r)$ with $D_{q,t}$.
\end{exercise}

\begin{definition}[DAHA Fourier assignment]
Define the DAHA Fourier assignment on the torus, winding, and Hecke
generators by
\[
\mathcal F(X_j)=Y_j^{-1},\qquad
\mathcal F(Y_j)=X_j,\qquad
\mathcal F(T_i)=T_i.
\]
For $\rho=(r-1,r-2,\ldots,0)$ and a partition $\lambda$, put
\[
\xi_\lambda=(q^{\lambda_1}t^{r-1},
q^{\lambda_2}t^{r-2},\ldots,q^{\lambda_r}).
\]
\end{definition}

\begin{exercise}[Macdonald Fourier duality]
Verify the DAHA relations after applying $\mathcal F$, prove that it
extends to an automorphism, and construct its
intertwining transform on the polynomial representation.  Normalize
$P_\lambda$ by its value at $t^\rho$ and prove the evaluation symmetry
\[
\frac{P_\lambda(\xi_\mu;q,t)}{P_\lambda(t^\rho;q,t)}
=
\frac{P_\mu(\xi_\lambda;q,t)}{P_\mu(t^\rho;q,t)}.
\]
Derive the duality between multiplication recurrences in $\lambda$ and
the commuting Macdonald difference equations in the spectral variable.
\end{exercise}

\begin{exercise}[Macdonald spectral theory]
Order monomial symmetric functions by dominance and prove that $D_{q,t}$
is triangular with the displayed diagonal entries.  For generic $q,t$,
deduce existence and uniqueness of $P_\lambda$.  Prove directly from the
torus integral that $D_{q,t}$ is symmetric when $0<q,t<1$, and deduce
orthogonality and joint diagonalization for symmetric Laurent polynomials
in the commuting $Y_i$.
\end{exercise}

\begin{definition}[Symmetric-function Macdonald pairing]
Let $\Lambda_{\rm sym}=\C[p_1,p_2,\ldots]$ be graded by
$\deg p_r=r$.  For a partition $\mu$, let $m_r(\mu)$ be the number of its
parts equal to $r$, and define
\[
z_\mu=\prod_{r\geq1}r^{m_r(\mu)}m_r(\mu)!,
\qquad
p_\mu=\prod_jp_{\mu_j},
\qquad
p_r(x)=\sum_i x_i^r.
\]
Define the symmetric-function Macdonald pairing by
\[
\langle p_\mu,p_\nu\rangle^{\rm sym}_{q,t}
=\delta_{\mu\nu}z_\mu
\prod_j\frac{1-q^{\mu_j}}{1-t^{\mu_j}}.
\]
For a cell $s=(i,j)$ in the diagram of $\lambda$, define its arm and leg by
\[
a_\lambda(s)=\lambda_i-j,
\qquad
\ell_\lambda(s)=\lambda'_j-i,
\]
and put
\[
b_\lambda(q,t)
=\prod_{s\in\lambda}
\frac{1-q^{a_\lambda(s)}t^{\ell_\lambda(s)+1}}
{1-q^{a_\lambda(s)+1}t^{\ell_\lambda(s)}}.
\]
\end{definition}

\begin{exercise}[Stable Macdonald pairing and Cauchy identity]
Prove that the finite-variable Macdonald polynomials are stable under
adjoining a zero variable and hence determine elements $P_\lambda$ of
$\Lambda_{\rm sym}$.  Put $Q_\lambda=b_\lambda(q,t)P_\lambda$.
Use the adjointness of the DAHA difference operators
and triangularity to prove
\[
\langle P_\lambda,P_\mu\rangle^{\rm sym}_{q,t}
=\delta_{\lambda\mu}b_\lambda(q,t)^{-1}.
\]
Deduce that $P_\lambda$ and $Q_\lambda$ are dual.  Expand the reproducing
kernel in the power sums and use the $q$-binomial identity to prove
\[
\sum_\lambda P_\lambda(a;q,t)Q_\lambda(b;q,t)
=\prod_{i,j}\frac{(t a_i b_j;q)_\infty}
{(a_i b_j;q)_\infty}
\]
for finite alphabets with $|a_i b_j|<1$.
\end{exercise}

\begin{definition}[Finite Macdonald measures]
Let $0\leq q,t<1$, let $a=(a_1,\ldots,a_N)$ and
$b=(b_1,\ldots,b_M)$ be nonnegative with $a_i b_j<1$, and let
$P_\lambda,Q_\lambda$ be the preceding Macdonald functions.  Define
\[
\Pi(a;b)=\sum_\lambda
P_\lambda(a;q,t)Q_\lambda(b;q,t),
\qquad
\mathbb P_{q,t}(\lambda)
=\frac{P_\lambda(a;q,t)Q_\lambda(b;q,t)}{\Pi(a;b)}.
\]
\end{definition}

\begin{exercise}[One Macdonald observable]
Apply the first Macdonald difference operator in the $a$ variables to the
Cauchy identity.  For distinct $a_i$, prove
\[
\mathbb E_{q,t}\!\left[
\sum_{i=1}^Nq^{\lambda_i}t^{N-i}\right]
=\sum_{i=1}^N
\prod_{j\ne i}\frac{t a_i-a_j}{a_i-a_j}
\prod_{k=1}^M\frac{1-a_i b_k}{1-t a_i b_k},
\]
and extend the identity to coincident parameters by continuity.  Verify
the Schur specialization $q=t$ and the $q$-Whittaker specialization $t=0$.
For samples from the finite measure, identify the displayed expectation as
a directly estimable statistic and state the positivity, truncation, and
sampling assumptions required for comparison with the formula.
\end{exercise}

\begin{definition}[Multiplicative theta function]
For $|p|<1$, define
\[
\theta(z;p)=(z;p)_\infty(p/z;p)_\infty.
\]
\end{definition}

\begin{definition}[Elliptic Ruijsenaars operator]
For type $A_{r-1}$, define
\[
D_{p,q,t}^{\mathrm{ell}}
=\sum_{i=1}^r
\left(
\prod_{j\neq i}
\frac{\theta(t x_i/x_j;p)}{\theta(x_i/x_j;p)}
\right)T_{q,x_i}.
\]
\end{definition}

\begin{exercise}[Ruijsenaars degeneration]
Prove directly from $\theta(z;0)=1-z$ that
\[
D_{p,q,t}^{\mathrm{ell}}
\longrightarrow D_{q,t}
\qquad(p\to0).
\]
Set $q=e^{-\varepsilon}$ and $t=e^{-g\varepsilon}$ and derive the
Heckman--Opdam differential operator from the second-order term as
$\varepsilon\to0$.  Take the rational scaling limit of its hyperbolic
coefficients and obtain the rational Calogero--Moser operator.
\end{exercise}

\begin{exercise}[The Calogero--Sutherland limit]
Write $x_i=e^{z_i}$, set $q=e^{-\epsilon}$ and $t=e^{-g\epsilon}$, and expand
$D_{q,t}$ through order $\epsilon^2$.  After subtracting its scalar value
on $1$ and removing the center-of-mass drift, obtain
\[
\mathcal L_g=
\sum_i\partial_{z_i}^2
+g\sum_{i<j}\coth\!\left(\frac{z_i-z_j}{2}\right)
 (\partial_{z_i}-\partial_{z_j}).
\]
On the circle, put $z_i=i\theta_i$ and conjugate by
\[
\Psi_0(\theta)=
\prod_{i<j}\left|2\sin\frac{\theta_i-\theta_j}{2}\right|^g.
\]
Recover, up to its ground-state energy, the Sutherland Hamiltonian
\[
H_g=-\sum_i\partial_{\theta_i}^2
+\frac{g(g-1)}2\sum_{i<j}
 \frac1{\sin^2((\theta_i-\theta_j)/2)}.
\]
Compute the eigenfunctions of degrees zero and one and, for $r=2$, degree
two.  Identify their Sutherland eigenvalues.
\end{exercise}

\begin{exercise}[Hypergeometric degeneration]
Set $q=e^{-\varepsilon}$.  Prove
\[
\lim_{\varepsilon\to0^+}
\frac{(q^a;q)_m}{(1-q)^m}=(a)_m,
\qquad
\lim_{q\to1^-}
(1-q)^{1-a}\frac{(q;q)_\infty}{(q^a;q)_\infty}
=\Gamma(a).
\]
Deduce
\[
\lim_{q\to1^-}
{}_r\phi_{r-1}\!\left(
\begin{matrix}q^{a_1},\ldots,q^{a_r}\\
q^{b_1},\ldots,q^{b_{r-1}}
\end{matrix};q,z\right)
=
{}_rF_{r-1}\!\left(
\begin{matrix}a_1,\ldots,a_r\\
b_1,\ldots,b_{r-1}
\end{matrix};z\right).
\]
Derive Gauss's summation from the $q$-Gauss summation.  Prove
the corresponding degeneration of the basic hypergeometric $q$-difference
equation to the Gauss equation.
\end{exercise}

\begin{definition}[\texorpdfstring{$BC_r$}{BCr} root-system weight]
Write $f(z^{\pm m})=f(z^m)f(z^{-m})$ and
\[
f(z_i^{\pm1}z_j^{\pm1})
=\prod_{\epsilon,\delta=\pm1}f(z_i^\epsilon z_j^\delta).
\]
For $|q|,|t|,|a_s|<1$, define
\[
\begin{aligned}
\Delta_{BC_r}(z)={}&
\prod_{i=1}^r
 \frac{(z_i^{\pm2};q)_\infty}
      {\prod_{s=1}^4(a_sz_i^{\pm1};q)_\infty}\\
&\times
\prod_{1\leq i<j\leq r}
 \frac{(z_i^{\pm1}z_j^{\pm1};q)_\infty}
      {(t z_i^{\pm1}z_j^{\pm1};q)_\infty}.
\end{aligned}
\]
Use the normalized torus measure
\[
\frac1{2^rr!}\prod_{i=1}^r\frac{dz_i}{2\pi iz_i}.
\]
\end{definition}

\begin{definition}[Koornwinder operator]
For $z=(z_1,\ldots,z_r)$ put
\[
A_i(z)=
\frac{\prod_{s=1}^4(1-a_sz_i)}
{(1-z_i^2)(1-qz_i^2)}
\prod_{j\ne i}
\frac{(1-tz_iz_j)(1-tz_i/z_j)}
{(1-z_iz_j)(1-z_i/z_j)}.
\]
On Laurent polynomials invariant under permutations and inversions define
\[
\mathcal D_K
=\sum_{i=1}^r\left[
A_i(z)(T_{q,z_i}-1)
+A_i(z^{-1})(T_{q,z_i}^{-1}-1)\right].
\]
The \emph{Koornwinder polynomial} $K_\lambda$ is the Weyl-invariant
Laurent polynomial triangular in orbit sums and satisfying
\[
\mathcal D_KK_\lambda
=\left(\sum_{i=1}^r
\left[
a_1a_2a_3a_4q^{-1}t^{2r-i-1}(q^{\lambda_i}-1)
+t^{i-1}(q^{-\lambda_i}-1)
\right]
\right)K_\lambda,
\]
with the scalar eigenvalue fixed by the displayed convention.
\end{definition}

\begin{exercise}[\texorpdfstring{$BC_r$}{BCr} weight]
Prove that $\Delta_{BC_r}$ is invariant under permutations and independent
inversions of the $z_i$.  Identify its factors with the roots
$\pm e_i$, $\pm2e_i$, and $\pm e_i\pm e_j$.  For $r=1$, recover the
Askey--Wilson weight and its order-two Weyl normalization, and identify
$\mathcal D_K$ with $\mathcal D_{AW}$.
\end{exercise}

\begin{exercise}[Koornwinder spectral theory]
Prove that $\mathcal D_K$ is
symmetric for the $\Delta_{BC_r}$ inner product and that its triangular
eigenfunctions are mutually orthogonal when the displayed eigenvalues are
distinct.
\end{exercise}

\begin{exercise}[\texorpdfstring{$BC_r$}{BCr} limit]
Set $t=q^g$
with $g\in\N$, reduce the quotients of $q$-products to finite products,
and compute the limit $q\to1^-$ with $z_i=e^{i\theta_i}$.
\end{exercise}

\begin{exercise}[Askey--Wilson spectrum]
Fix $M\geq4$ and parameters $a,b,c,d\in(-1,1)$ for which the
Askey--Wilson weight is positive and the first $M+1$ eigenvalues are
distinct.  Let $\mathcal H_M$ be the span of the normalized Askey--Wilson
polynomials $p_0,\ldots,p_M$ with their positive Askey--Wilson inner
product.  Model a finite programmable $q$-difference
simulator by the self-adjoint Hamiltonian
$\widehat H_M=\varepsilon\mathcal D_{AW}|_{\mathcal H_M}$, where
$\varepsilon>0$ is a calibrated energy scale.  Prove that its spectral
values are
\[
E_n=\varepsilon(q^{-n}-1)(1-abcdq^{n-1}).
\]
Write the three-term recurrence in the form
\[
(z+z^{-1})p_n=A_np_{n+1}+B_np_n+C_np_{n-1}
\]
and compute the corresponding normalized matrix elements.  Deduce that a
periodic perturbation proportional to the compression of $z+z^{-1}$ to
$\mathcal H_M$ produces only neighboring transitions, with resonance frequencies
$|E_{n+1}-E_n|/\hbar$ and intensities determined by $A_n$ and $C_{n+1}$.
Assuming $\varepsilon$ is known and at least three nondegenerate lines are
resolved, derive the equations satisfied by $q$ and the product $abcd$.
Compute the Jacobian of two independent gap ratios and state the
nonvanishing and parameter-range hypotheses under which these two
quantities are locally identifiable.

Form the data map from $(a,b,c,d)$ to four selected squared normalized
neighboring matrix elements.  Prove that it is invariant under permutation
of the parameters.  At a chosen parameter point, compute its differential
in the four elementary symmetric functions and state the full-rank
hypothesis under which the measured strengths locally determine the
unordered set $\{a,b,c,d\}$.  When this rank condition fails, identify the
remaining fiber instead of asserting identifiability; state also that
squared strengths cannot distinguish any additional discrete ambiguity in
that fiber.  State that these are predictions for the engineered simulator.
\end{exercise}

\begin{exercise}[Macdonald--Ruijsenaars spectrum]
Let $0<q,t<1$, fix $M\geq0$, and define
\[
\mathcal H_{q,t}^{(M)}
=\operatorname{span}\{P_\lambda(x;q,t):|\lambda|\leq M\}
\]
with the positive Macdonald inner product.  Put
\[
d_0=\sum_{i=1}^r t^{r-i},
\qquad
\widehat H_{q,t}^{(M)}
=\varepsilon(d_0-D_{q,t})|_{\mathcal H_{q,t}^{(M)}},
\qquad \varepsilon>0.
\]
Identify $D_{q,t}$ with the first symmetric function of the commuting
DAHA winding operators and derive its compatibility with the Hecke
exchange relations.  Prove that $\mathcal H_{q,t}^{(M)}$ is invariant,
that $\widehat H_{q,t}^{(M)}$ is self-adjoint, and
that the normalized Macdonald states have energies
\[
E_\lambda
=\varepsilon\sum_{i=1}^r
t^{r-i}(1-q^{\lambda_i}).
\]
\end{exercise}

\begin{exercise}[Macdonald transition spectrum]
Model a finite programmable $q$-difference simulator by this operator.
Use the Macdonald Pieri identity for
$\rho_1=x_1+\cdots+x_r$ and its adjoint to determine every nonzero matrix
element of the compression of $\rho_1+\rho_1^*$ to
$\mathcal H_{q,t}^{(M)}$.  For a periodic drive proportional to this
compressed observable, derive the resonance frequencies
\[
\Omega_{\mu\lambda}=|E_\mu-E_\lambda|/\hbar
\]
and the relative intensities as squared normalized Macdonald Pieri
coefficients.  Explain that recovery of
$q,t$ requires known $\varepsilon$ and $r$, resolved partition labels,
sufficiently many nondegenerate transitions, and a fixed convention for
any discrete parameter ambiguity, and that the prediction concerns the
engineered simulator.
\end{exercise}

\begin{exercise}[Macdonald--Jack degeneration]
Set $q=e^{-\delta}$ and $t=e^{-g\delta}$; after the scalar
and center-of-mass normalizations in the Calogero--Sutherland limit
exercise, prove that
the rescaled spectral gaps and transition coefficients converge to the
Sutherland energies and Jack Pieri coefficients.  State how measured line
locations and strengths test this degeneration.  Explain that recovery of
the deformation parameter requires resolved partition labels and
sufficiently many nondegenerate transitions.
\end{exercise}

\begin{definition}[Elliptic shifted factorials and gamma function]
For $|p|,|q|<1$, define
\[
\theta(z_1,\ldots,z_s;p)=\prod_{j=1}^s\theta(z_j;p),
\]
\[
(a;q,p)_m=\prod_{j=0}^{m-1}\theta(aq^j;p),
\]
\[
(a_1,\ldots,a_s;q,p)_m
=\prod_{i=1}^s(a_i;q,p)_m,
\]
and
\[
\Gamma(z;p,q)=\prod_{j,k\geq0}
\frac{1-p^{j+1}q^{k+1}/z}{1-p^jq^kz}.
\]
\end{definition}

\begin{definition}[Deformation hierarchy]
For a series $\sum_m c_m$, set $h_m=c_{m+1}/c_m$.  It is
\emph{ordinary hypergeometric} when $h_m$ is rational in $m$, and
\emph{basic hypergeometric} when $h_m$ is rational in $q^m$.  It is
\emph{elliptic hypergeometric} when, after writing $x=q^m$, its term ratio
has the form
\[
h(x)=z\frac{\prod_{i=1}^s\theta(a_ix;p)}
{\prod_{i=1}^s\theta(b_ix;p)}.
\]
Its parameters are \emph{balanced} when
\[
\prod_{i=1}^s a_i=\prod_{i=1}^s b_i.
\]
By theta quasi-periodicity this is equivalent to
\[
h(px)=h(x).
\]
\end{definition}

\begin{exercise}[Elliptic degenerations]
Prove
\[
\theta(pz;p)=-z^{-1}\theta(z;p),
\qquad
\Gamma(qz;p,q)=\theta(z;p)\Gamma(z;p,q),
\]
\[
\Gamma(pz;p,q)=\theta(z;q)\Gamma(z;p,q),
\qquad
\Gamma(z;p,q)\Gamma(pq/z;p,q)=1.
\]
Prove
\[
\theta(z;0)=1-z,
\qquad
(a;q,0)_m=(a;q)_m,
\qquad
\Gamma(z;0,q)=\frac1{(z;q)_\infty}.
\]
\end{exercise}

\begin{definition}[Frenkel--Turaev sum]
For $a^2q^{n+1}=bcde$, define
\[
\begin{aligned}
{}_{10}V_9(a;b,c,d,e,q^{-n};q,p)
=\sum_{m=0}^n{}&
\frac{\theta(aq^{2m};p)}{\theta(a;p)}\\
&\times
\frac{(a,b,c,d,e,q^{-n};q,p)_m}
{(q,aq/b,aq/c,aq/d,aq/e,aq^{n+1};q,p)_m}q^m.
\end{aligned}
\]
\end{definition}

\begin{exercise}[Theta addition and Frenkel--Turaev summation]
Prove the multiplicative theta addition formula
\[
\begin{aligned}
&\theta(xz,x/z,yw,y/w;p)-\theta(xw,x/w,yz,y/z;p)\\
&\hspace{4em}=\frac{y}{z}\theta(xy,x/y,zw,z/w;p)
\end{aligned}
\]
by comparing quasi-periodicity and zeros.  Use it to telescope the
difference between the $n$th and $(n-1)$st terminating sums and prove
\[
{}_{10}V_9(a;b,c,d,e,q^{-n};q,p)
=\frac{(aq,aq/(bc),aq/(bd),aq/(cd);q,p)_n}
{(aq/b,aq/c,aq/d,aq/(bcd);q,p)_n}.
\]
Take $p\to0$ and identify the terminating very-well-poised basic
hypergeometric summation.
\end{exercise}

\begin{definition}[Elliptic beta kernel]
Write
\[
\Gamma(tz^{\pm1};p,q)=\Gamma(tz;p,q)\Gamma(t/z;p,q),
\qquad
\kappa_{p,q}=\frac{(p;p)_\infty(q;q)_\infty}{2}.
\]
For $|t_j|<1$ with $\prod_{j=1}^6t_j=pq$, define
\[
\mathcal B(t_1,\ldots,t_6)
=\kappa_{p,q}\int_{\T}
\frac{\prod_{j=1}^6\Gamma(t_jz^{\pm1};p,q)}
{\Gamma(z^{\pm2};p,q)}\,\frac{dz}{2\pi iz}.
\]
\end{definition}

\begin{exercise}[Elliptic beta recurrence]
Use the elliptic-gamma difference equations to derive the parameter-shift
recurrences of $\mathcal B$.
\end{exercise}

\begin{exercise}[Elliptic beta integral]
Evaluate a residue degeneration for which
two parameters multiply to $pq$ and use analytic continuation to prove
\[
\mathcal B(t_1,\ldots,t_6)
=\prod_{1\leq i<j\leq6}\Gamma(t_it_j;p,q).
\]
Take $p\to0$ and identify the resulting basic hypergeometric beta
integral.
\end{exercise}

\begin{definition}[Elliptic integral operators]
For $|t|<1$ define
\[
(\mathcal M_t f)(w)
=\frac{\kappa_{p,q}}{\Gamma(t^2;p,q)}
\int_{\T}
\frac{\Gamma(tw^{\pm1}z^{\pm1};p,q)}
{\Gamma(z^{\pm2};p,q)}f(z)\,\frac{dz}{2\pi iz},
\]
and
\[
\mathcal D_s(y,w)
=\Gamma\!\left((pq)^{1/2}s^{-1}y^{\pm1}w^{\pm1};p,q\right).
\]
\end{definition}

\begin{exercise}[Elliptic star--triangle relation]
Apply the elliptic beta integral to the internal integration variable and
prove the operator identity
\[
\mathcal M_s\mathcal D_{st}(y,\cdot)\mathcal M_t
=\mathcal D_t(y,\cdot)\mathcal M_{st}\mathcal D_s(y,\cdot)
\]
whenever the contours separate reciprocal pole sequences and the
parameters satisfy the corresponding modulus bounds.  Interpret the
kernels as edge weights and $\mathcal D$ as a site weight of a continuous
spin model.  Prove that the two local triangular partition functions are
equal and that the row transfer operators obtained by the two
factorizations commute.  For a finite periodic strip, deduce equality of
the transfer spectra and partition functions under a star--triangle move;
state the positivity regime and contour assumptions under which these are
statistical-mechanical observables.
\end{exercise}
